\documentclass[UTF8,zihao=-4,openany,oneside,fontset=windows]{ctexbook}
\usepackage[a4paper,top=24mm,bottom=24mm,left=24mm,right=22mm]{geometry}
\usepackage{amsmath,amssymb,bm,mathtools}
\usepackage{graphicx,booktabs,longtable,array,enumitem}
\usepackage{xcolor,fancyhdr,fvextra}
\usepackage{xurl}
\usepackage[unicode,colorlinks=true,linkcolor=blue!45!black,urlcolor=blue!55!black]{hyperref}
\usepackage{bookmark}
\setmonofont{Consolas}
\definecolor{inkblue}{HTML}{164766}
\ctexset{chapter={format=\Huge\bfseries\color{inkblue}},section={format=\Large\bfseries\color{inkblue}},subsection={format=\large\bfseries}}
\setlength{\parskip}{0.35em}
\setlength{\headheight}{15pt}
\pagestyle{fancy}
\fancyhf{}
\fancyhead[L]{\small 国赛 A 题通用建模手册}
\fancyhead[R]{\small 零基础详解版}
\fancyfoot[C]{\thepage}
\setlist{nosep,leftmargin=2em}
\allowdisplaybreaks
\emergencystretch=3em
\newcommand{\dd}{\mathrm d}
\newcommand{\R}{\mathbb R}
\hypersetup{pdftitle={数学建模国赛A题通用公式与代码手册},pdfauthor={Codex},pdfsubject={零基础公式、参数、条件与逐行代码解释}}

\begin{document}
\begin{titlepage}
\centering
\vspace*{22mm}
{\large\color{inkblue} 数学建模 · 公式与计算工作手册\par}
\vspace{12mm}
{\Huge\bfseries 国赛 A 题\par}
\vspace{5mm}
{\Huge\bfseries 通用公式与代码\par}
\vspace{10mm}
{\LARGE 零基础详解版\par}
\vspace{16mm}
{\large ODE 求解\quad 一维热传导有限差分\quad 非线性优化\par}
\vspace{4mm}
{\large 力学\quad 光学\quad 热学\quad 流体\quad 概率统计\par}
\vspace{4mm}
{\large 基础定理\quad 代表性前沿文献\quad Python 逐行注释\par}
\vfill
\begin{minipage}{0.82\textwidth}
\large 每个公式都从“它在说什么”开始，说明每个符号、单位、使用条件与常见误用。先理解物理，再选择算法，最后验证结果。
\end{minipage}
\vfill
{\large 2026 年 9 月 10 日\par}
\vspace{5mm}
{\small 通用教学参考；不针对某届题目，不承诺涵盖全部学科公式。\par}
\end{titlepage}
\frontmatter
\tableofcontents
\mainmatter
\chapter{先学会读公式：从题目到模型}
\section{这本手册怎么用}
这是面向全国大学生数学建模竞赛本科 A 题风格的通用参考，不是某一年题目的答案，也不是涵盖所有学科、所有文献的公式全集。题目可能横跨多个领域，应由题目事实决定模型。优先学守恒定律、ODE、热方程、参数拟合和误差检查，再按题目补充物理与统计。前沿方法是选读，不能代替可靠的基础模型。

每个公式按“意思—符号—条件—例子或注意”来读。字母在不同章节可能有不同含义，以公式附近的定义为准。代码中的数字都是教学示例，不能当作题目提供的数据或材料常数。

推荐顺序：第一遍看本章与对应的物理小节；第二遍运行 ODE 和热传导模板；第三遍学习优化与不确定性。每完成一步，先核对单位和一个简单算例，再增加复杂度。

\section{最常遇到的符号}
\begin{description}[style=nextline]
\item[$x,t,T$] 常分别表示位置、时间、温度。它们是变量：会随情景改变。
\item[$a,k,\theta$] 常表示常数或待估参数。参数在一次模拟内固定，但不同材料或拟合结果可能不同。
\item[$f(x)$] 输入一个 $x$，按规则 $f$ 算出一个结果。例如温度可以是位置的函数。
\item[$\mathrm{d}y/\mathrm{d}t$] $y$ 对时间的变化速度；不是 $y/t$。若 $y$ 用米、$t$ 用秒，导数用米每秒。
\item[$\partial T/\partial t$] 多变量函数中的偏导：固定位置等其他变量，只看时间变化。
\item[$\sum_{i=1}^{n}a_i$] 把编号从 1 到 $n$ 的数相加；$i$ 是编号，$n$ 是项数。
\item[$\int_a^b f(x)\,\mathrm{d}x$] 从 $a$ 到 $b$ 累加连续的小贡献；积分的单位是 $f$ 的单位乘 $x$ 的单位。
\item[$\boldsymbol{x},\|\boldsymbol{x}\|_2$] 粗体是多个数组成的向量；下标 2 的范数是各分量平方和开根号。
\item[$\nabla f,\nabla\cdot\boldsymbol{v},\nabla^2 f$] 依次为梯度（各方向变化率）、散度（局部净流出程度）、拉普拉斯算子（各方向二阶偏导之和）。在一维，拉普拉斯就是二阶导数。
\item[$\approx,\propto,\in,\forall$] 依次为近似相等、成比例、属于、对所有。
\item[$\arg\min$] 返回使目标最小的输入；$\min$ 返回最小的目标值，两者不同。
\item[$\exp(z),\ln z$] 以自然常数 $e$ 为底的指数和对数；指数的输入必须无量纲，对数输入须为正且无量纲。
\end{description}

\section{单位与量纲：最便宜的排错手段}
统一采用 SI：长度 m、时间 s、质量 kg、温度 K、力 N、能量 J、功率 W、压力 Pa。摄氏温差与开尔文温差数值相同，但绝对温度不同。参考 NIST 的 SI 说明：\url{https://www.nist.gov/pml/owm/metric-si/si-units}。
\subsection{温度换算}
\begin{equation}T_{\mathrm K}=T_{\mathrm C}+273.15.\end{equation}
这里 $T_{\mathrm K}$ 是开尔文温度的数值，$T_{\mathrm C}$ 是摄氏温度的数值，273.15 是两种温标零点的差。例：$20\,^{\circ}\mathrm C$ 是 $293.15\,\mathrm K$。辐射和理想气体公式必须用 K。
\subsection{量纲齐次}
\begin{equation}[a+b]=[a]=[b],\qquad [ab]=[a][b],\qquad [a/b]=[a]/[b].\end{equation}
$[a]$ 表示量 $a$ 的物理量纲，$a,b$ 是物理量。只有相同量纲能相加；除法要求分母非零。例如速度乘时间是长度，不能把温度直接加到长度上。无量纲不等于数值一定很小。

\section{所有物理模型的主线：守恒}
\begin{equation}\frac{\mathrm d S}{\mathrm dt}=Q_{\rm in}-Q_{\rm out}+Q_{\rm gen}.\end{equation}
$S$ 是系统储存量，可以是质量 kg 或能量 J；$t$ 是时间 s；$Q_{\rm in},Q_{\rm out},Q_{\rm gen}$ 分别是流入、流出、内部净生成速率，单位为 $S$ 的单位每秒。意思是“存量增加=进来−出去+内部产生”。必须先固定系统边界。总能量守恒时，化学能转成热不是凭空产生总能量，而是所追踪的热能增加。
例：水箱每秒进 3 kg、出 2 kg，没有内部生成，则每秒多 1 kg。

\section{泰勒展开、数值微分和积分}
\subsection{局部线性近似与泰勒公式}
\begin{equation}f(x+h)=f(x)+hf'(x)+\frac{h^2}{2}f''(\xi).\end{equation}
$f$ 是至少二阶连续可微的标量函数，$x$ 是展开点，$h$ 是输入改变量，$f',f''$ 是一阶和二阶导数，$\xi$ 是 $x$ 与 $x+h$ 之间某点。最后一项是精确余项而非已知常数。丢掉余项得到局部直线近似；$h$ 不小或曲率很大时误差会大。所有项单位都与 $f$ 相同。
\subsection{一阶导数中心差分}
\begin{equation}f'(x)\approx\frac{f(x+h)-f(x-h)}{2h}.\end{equation}
$h>0$ 是步长，$x$ 是求导点，$f$ 是函数；导数单位是输出单位除输入单位。函数足够光滑时截断误差为二阶，即步长减半通常使这部分误差约变为四分之一。噪声数据不能无限减小 $h$，否则差值除小数会放大噪声。
\subsection{复合梯形积分}
\begin{equation}\int_a^b f(x)\,\mathrm dx\approx h\left[\frac{f(x_0)+f(x_N)}2+\sum_{i=1}^{N-1}f(x_i)\right],\quad h=\frac{b-a}{N},\quad x_i=a+ih.\end{equation}
$a,b$ 是积分下上限且 $b>a$，$N$ 是正整数区间数，$h$ 是等距间隔，$x_i$ 是第 $i$ 个采样点，$f(x_i)$ 是该点函数值。两端只算半份，中间算整份。例：对恒定 10 W 功率积 3 s，能量为 30 J。非等距数据必须分别用各段宽度。

\section{向量、矩阵与线性方程}
\subsection{点积与范数}
\begin{equation}\boldsymbol a^\mathsf T\boldsymbol b=\sum_{i=1}^n a_i b_i,\qquad \|\boldsymbol a\|_2=\sqrt{\sum_{i=1}^n a_i^2}.\end{equation}
$\boldsymbol a,\boldsymbol b$ 是具有 $n$ 个实数分量的向量，$a_i,b_i$ 是第 $i$ 个分量，$\mathsf T$ 表示转置。点积用来计算投影与做功；范数类似多维长度。直接把米、秒、开尔文混合进范数没有合理的物理含义，应先缩放。
\subsection{线性系统}
\begin{equation}A\boldsymbol x=\boldsymbol b.\end{equation}
$A$ 是 $m\times n$ 系数矩阵，$\boldsymbol x$ 是 $n$ 维未知向量，$\boldsymbol b$ 是 $m$ 维已知向量。$m=n$ 且 $A$ 可逆时解唯一；其他情形需依据矩阵的秩与方程一致性判断，也可能有唯一解。代码应用线性求解器，不要为了求解而显式计算逆矩阵。热传导隐式离散就会产生这种系统。
\subsection{奇异值分解与条件数}
\begin{equation}A=U\Sigma V^\mathsf T,\qquad \kappa_2(A)=\frac{\sigma_{\max}}{\sigma_{\min}}.\end{equation}
$A$ 是实矩阵，$U,V$ 为正交矩阵，$\Sigma$ 为非负奇异值组成的对角形矩阵；$\sigma_{\max},\sigma_{\min}$ 是最大与最小奇异值，条件数这里针对满秩方阵。$\kappa_2$ 无量纲，越大越容易把输入误差放大；最小奇异值为零时奇异。SVD 也可用于降维，但不能因方便而丢掉关键小尺度信号。

\section{无量纲化与相似}
\subsection{缩放变量}
\begin{equation}x^*=\frac{x}{L},\qquad t^*=\frac{t}{\tau},\qquad T^*=\frac{T-T_{\rm ref}}{\Delta T}.\end{equation}
$x$ 为位置 m，$L>0$ 为特征长度 m；$t$ 为时间 s，$\tau>0$ 为特征时间 s；$T,T_{\rm ref}$ 为温度，$\Delta T\ne0$ 是选定温差。星号表示无量纲变量，不表示最优值。选择典型尺度后，变量通常在数量级 1 附近，有利于求解。
\subsection{Buckingham $\Pi$ 定理}
\begin{equation}F(q_1,\ldots,q_n)=0\quad\Longrightarrow\quad G(\Pi_1,\ldots,\Pi_{n-r})=0.\end{equation}
$q_1,\ldots,q_n$ 是 $n$ 个物理量，$F$ 是量纲齐次关系；$r$ 是量纲指数矩阵的秩，$\Pi_j$ 为由原物理量幂乘构成的无量纲组合，$G$ 是重新表达后的关系。定理在可用这些量完整表述、量纲变换一致的前提下说明独立组合数，不能自动给出 $G$ 的具体函数。例：阻力研究常转成阻力系数与雷诺数关系。参考 MIT 量纲分析教学材料及物理教材，本式为标准定理表述。

\section{误差、收敛与灵敏度}
\subsection{绝对误差和稳健相对误差}
\begin{equation}e_{\rm abs}=|y_{\rm num}-y_{\rm ref}|,\quad e_{\rm rel}=\frac{|y_{\rm num}-y_{\rm ref}|}{\max(|y_{\rm ref}|,s_0)}.\end{equation}
$y_{\rm num}$ 是数值结果，$y_{\rm ref}$ 是解析解或高精度参考，$s_0>0$ 为有同单位的小参考尺度，用于避免除零；$e_{\rm abs}$ 与 $y$ 同单位，$e_{\rm rel}$ 无量纲。当参考接近零时应同时报告绝对误差和所用尺度，不能随意挑尺度美化结果。
\subsection{观测收敛阶}
\begin{equation}p_{\rm obs}=\log_2\!\left(\frac{|u_h-u_{h/2}|}{|u_{h/2}-u_{h/4}|}\right).\end{equation}
$u_h,u_{h/2},u_{h/4}$ 是同一物理位置和时刻、三种等比细化网格得到的同一标量输出；$h$ 是网格尺度，$p_{\rm obs}$ 是估计收敛阶，无量纲。需要误差处于单一主导幂次的渐近区间且差值非零。不能拿不同位置的输出相减；空间研究时要压低时间误差。
\subsection{归一化局部灵敏度}
\begin{equation}S_j=\frac{\theta_j}{y}\frac{\partial y}{\partial\theta_j}.\end{equation}
$y\ne0$ 是模型标量输出，$\theta_j$ 是第 $j$ 个参数，$S_j$ 无量纲，表示参数相对变化引起的输出相对变化的一阶比例。若 $S_j=2$，小幅增加参数 1\% 通常导致输出约增 2\%；这是局部近似。输出或参数接近零时应采用未归一化导数或预先指定的非零参考尺度。
\subsection{一阶不确定性传播}
\begin{equation}\operatorname{Var}(y)\approx J\Sigma_\theta J^\mathsf T,\qquad J_j=\frac{\partial f}{\partial\theta_j}.\end{equation}
$y=f(\boldsymbol\theta)$ 是参数到输出的函数，$J$ 是在参数均值处计算的行向量，$\Sigma_\theta$ 是参数协方差矩阵，$\operatorname{Var}(y)$ 是输出方差，单位是输出单位的平方。适合小扰动、局部近似线性情况；保留协方差才能反映参数相关性。强非线性可改用蒙特卡洛。

\section{从题面开始的七步}
\begin{enumerate}
\item 写清要预测什么、优化什么、评价什么，输出必须能回答题目。
\item 列符号表：已知量、待估参数、状态变量、控制变量、单位和来源。
\item 画系统边界，写守恒，再补材料关系和几何关系。
\item 写初始条件和边界条件；只写微分方程通常不够。
\item 用最简单特例和量纲检验公式，再代入真实数据。
\item 做参数估计、独立验证、网格与容差检查，报告误差和适用范围。
\item 优化前确认约束表达了真实限制，优化后检查可行性与鲁棒性。
\end{enumerate}

\chapter{常微分方程：从变化率走到随时间变化的结果}
\section{先弄懂状态、初值和变化率}
看到速度、温度、浓度随着时间变化，先问：系统现在的情况能用哪些数表示？这些数叫状态。状态随时间的变化规律叫常微分方程（ODE）。它不是直接告诉你答案，而是告诉你下一小段时间往哪个方向变化、变化多快。只有时间一个独立变量时是 ODE；温度同时随位置和时间变化时，一般要用偏微分方程（PDE）。

\subsection{一阶初值问题与方程组}
\begin{equation}
\frac{\mathrm d\boldsymbol y}{\mathrm dt}=\boldsymbol f(t,\boldsymbol y;\boldsymbol p),\qquad \boldsymbol y(t_0)=\boldsymbol y_0.
\end{equation}
这里 $t$ 是时间，常用单位 s；$t_0$ 是开始时间；$\boldsymbol y=(y_1,\ldots,y_m)^{\mathsf T}$ 是 $m$ 个状态组成的列向量，上标 $\mathsf T$ 表示转置；$m$ 是状态数，无量纲；$\boldsymbol y_0$ 是初始状态；$\boldsymbol p$ 是给定参数；$f_i$ 是第 $i$ 个状态的变化率，单位为 $y_i$ 的单位除以 s。不同状态可以有不同单位，例如一个状态是位移 m，另一个是速度 m/s。分号仅用于区分已知参数与状态，没有特殊运算意义。例：只有一个温度状态时，$y_1=T$，$f_1$ 的单位就是 K/s。

\subsection{二阶方程必须先拆成一阶}
\begin{equation}
m\ddot x+c\dot x+kx=F(t)\quad\Longrightarrow\quad
\frac{\mathrm d}{\mathrm dt}\begin{pmatrix}x\\v\end{pmatrix}
=\begin{pmatrix}v\\ {}[F(t)-cv-kx]/m\end{pmatrix},\quad v=\dot x.
\end{equation}
此处 $m$ 改为质量，单位 kg；$x$ 是位移 m；$v$ 是速度 m/s；上方一点、两点分别表示对时间求一阶、二阶导数；$c$ 是黏性阻尼系数 N$\cdot$s/m；$k$ 是弹簧刚度 N/m；$F$ 是外力 N。左右各项均为力。求解需要两个初值 $x(t_0)$ 和 $v(t_0)$。适用于质量集中、线性弹簧、小变形和线性阻尼的系统。例：$m=1$ kg、$c=0.2$ N$\cdot$s/m、$k=4$ N/m、$F=0$ 时，第一个返回量必须是速度，第二个返回量才是加速度。

\subsection{解析解是检查程序的尺子}
\begin{equation}
\dot T=-\beta(T-T_\infty),\qquad
T(t)=T_\infty+(T_0-T_\infty)e^{-\beta(t-t_0)}.
\end{equation}
$T$ 是物体温度，$T_\infty$ 是恒定环境温度，$T_0$ 是初始温度，均可用 $^\circ$C；$\beta>0$ 是冷却速率，单位 s$^{-1}$；$e$ 是自然常数，无量纲。指数里的 $\beta(t-t_0)$ 必须无量纲。适用于物体内部近似等温、环境恒定、换热参数恒定的集中参数模型。温差用摄氏度或开尔文数值相同；涉及辐射的四次方时必须改成 K。例：$T_0=90$、$T_\infty=20$、$\beta=0.08$，一段时间后温度向 20 渐近，不会在有限时间突然等于 20。

\begin{equation}
t_\star=t_0+\frac{1}{\beta}\ln\frac{T_0-T_\infty}{T_\star-T_\infty}.
\end{equation}
$T_\star$ 是目标温度，与前面温度同单位；$t_\star$ 是首次达到的时间 s；$\ln$ 是自然对数。仅在冷却且 $T_\infty<T_\star<T_0$ 时使用这个直接形式。例：上述参数下达到 40 摄氏度约需 15.66 s。目标为 10 摄氏度时不能把公式硬算下去，因为这个模型根本到不了。

\section{手工数值方法：把曲线切成许多小步}
\subsection{时间网格}
\begin{equation}
t_n=t_0+nh,\qquad n=0,1,\ldots,N,\qquad h=(t_f-t_0)/N.
\end{equation}
$t_f$ 是结束时间 s；$N$ 是总步数，无量纲正整数；$h$ 是时间步长 s；$n$ 是步编号；$\boldsymbol y_n$ 以后表示在 $t_n$ 上对真实状态 $\boldsymbol y(t_n)$ 的近似。下标 $n$ 不表示乘法。例：从 0 算到 10 s，用 100 步，就是每次前进 0.1 s。

\subsection{显式欧拉：照着当前斜率往前走}
\begin{equation}
\boldsymbol y_{n+1}=\boldsymbol y_n+h\boldsymbol f(t_n,\boldsymbol y_n).
\end{equation}
$\boldsymbol y_n$ 和 $\boldsymbol y_{n+1}$ 是旧、新状态；$h\boldsymbol f$ 是这段时间内的状态增量，与状态同单位；$\boldsymbol f$、$t_n$、$h$ 的含义同前。适合讲解与初步试算，精度不高。对于前述冷却例题，初始变化率为 $-5.6$ 摄氏度每秒，取 $h=0.1$ s，下一步就是 $90-0.56=89.44$ 摄氏度。每步更新只能使用旧状态；向量内各分量要一起更新。

\subsection{隐式欧拉：新状态也出现在右边}
\begin{equation}
\boldsymbol y_{n+1}=\boldsymbol y_n+h\boldsymbol f(t_{n+1},\boldsymbol y_{n+1}).
\end{equation}
符号与显式欧拉相同，但右侧用未知的新状态。每步要解线性或非线性方程，因此算一步更贵，却更适合快速衰减导致的刚性问题。例：$\dot y=-a y$ 时，其中 $a>0$ 的单位 s$^{-1}$，可直接整理为 $y_{n+1}=y_n/(1+ah)$。分母始终大于 1，因此不会出现显式欧拉在大步长下的放大振荡。

\subsection{经典四阶 Runge--Kutta：每步探路四次}
\begin{align}
\boldsymbol k_1&=\boldsymbol f(t_n,\boldsymbol y_n),\\
\boldsymbol k_2&=\boldsymbol f(t_n+h/2,\boldsymbol y_n+h\boldsymbol k_1/2),\\
\boldsymbol k_3&=\boldsymbol f(t_n+h/2,\boldsymbol y_n+h\boldsymbol k_2/2),\\
\boldsymbol k_4&=\boldsymbol f(t_n+h,\boldsymbol y_n+h\boldsymbol k_3),\\
\boldsymbol y_{n+1}&=\boldsymbol y_n+\frac h6(\boldsymbol k_1+2\boldsymbol k_2+2\boldsymbol k_3+\boldsymbol k_4).
\end{align}
$\boldsymbol k_1$ 至 $\boldsymbol k_4$ 是四次估计的斜率，各分量单位为相应状态单位每秒；其余符号同前。这里的 $k_i$ 不是弹簧刚度。前两次看起点和中点，第三次改进中点，第四次看终点，再加权平均。适用于解比较光滑、非刚性的问题。不要把这个固定步长 RK4 与软件中的自适应 RK45 当成完全相同的算法。

\subsection{误差阶：步长变小为什么有用}
\begin{equation}
E(h)=\max_{0\leq n\leq N}|y_n-y(t_n)|\approx C h^p,
\quad \frac{E(h)}{E(h/2)}\approx 2^p.
\end{equation}
$E$ 是某一个状态的最大绝对误差，单位同该状态；$C$ 是依赖方程和区间的常数，单位为状态单位除以 s$^p$；$p$ 是全局收敛阶，无量纲；$|\cdot|$ 是绝对值。解足够光滑、方法稳定、舍入误差尚未主导时，欧拉法 $p=1$，经典 RK4 为 $p=4$。例：RK4 步长减半，误差理想情况下约缩小 16 倍。这不是任意复杂模型都无条件满足的保证。各分量单位不同时，不应直接把它们的绝对误差取共同最大值，应该先归一化。

\begin{equation}
p_{\mathrm{obs}}=\frac{\ln(E(h)/E(h/2))}{\ln 2},\qquad
E_{\mathrm{fine}}\approx\frac{|y_{h/2}(t_f)-y_h(t_f)|}{2^p-1}.
\end{equation}
$p_{\mathrm{obs}}$ 是观察到的误差阶；$y_h$、$y_{h/2}$ 是两次计算在同一终点的值；$E_{\mathrm{fine}}$ 是细步长终点误差估计，单位同 $y$。第二式是主导误差项相同且已进入渐近区间时的步长加倍估计，不能把它当作严格误差上界。没有解析解时，至少比较三套步长，并确认误差比例趋于稳定。

\subsection{稳定性与刚性：细步长有时是被迫的}
\begin{equation}
\dot y=\lambda y,\quad y_{n+1}=(1+h\lambda)y_n,
\quad |1+h\lambda|\leq1.
\end{equation}
$\lambda$ 是线性增长或衰减系数，单位 s$^{-1}$，可以是复数；$h\lambda$ 无量纲；绝对值对复数表示模。最后是不让显式欧拉放大衰减模式的稳定条件。对实数 $\lambda<0$，它成为 $0<h\leq 2/|\lambda|$；等号处可能不衰减，因此实践中留余量。例：$\lambda=-1000$ s$^{-1}$ 时，哪怕只关心秒级慢变化，显式欧拉仍需不超过 0.002 s 的步长。

\begin{equation}
\kappa_{\mathrm{time}}=\frac{\max_j |\operatorname{Re}\lambda_j|}{\min_j |\operatorname{Re}\lambda_j|}.
\end{equation}
$\lambda_j$ 是局部雅可比矩阵 $\partial\boldsymbol f/\partial\boldsymbol y$ 的衰减特征值，单位 s$^{-1}$；$\operatorname{Re}$ 表示实部；$j$ 是模式编号；$\kappa_{\mathrm{time}}$ 无量纲。只对实部严格为负且不接近零的一组模式这样计算，它是时间尺度分离的提示，不是刚性的完整判据。存在正特征值说明可能真的不稳定；存在零特征值时分母不能照抄。刚性问题优先尝试 Radau 或 BDF，并比较解与计算成本。

\section{自动步长、事件和结果检查}
\begin{equation}
s_i=\mathrm{atol}_i+\mathrm{rtol}|y_i|,\qquad
\epsilon_i=\frac{e_i}{s_i}.
\end{equation}
$s_i$ 是第 $i$ 个状态的误差尺度，与 $y_i$ 同单位；$\mathrm{atol}_i$ 是绝对误差容限，同单位；$\mathrm{rtol}$ 是相对容限，无量纲；$e_i$ 是求解器估计的局部误差；$\epsilon_i$ 是归一化误差，无量纲。求解器根据整体归一化误差调节步长，这不等于承诺最终全局误差恰好小于这两个参数。例：位移为 $10^{-6}$ m 时，若把绝对容限设为 $10^{-3}$ m，设置很小的相对容限也无法挽救对微小位移的粗糙控制。

\begin{equation}
g(t_\star,\boldsymbol y(t_\star))=0.
\end{equation}
$g$ 是用户定义的事件函数，单位由物理定义决定；$t_\star$ 是事件发生时刻 s；$\boldsymbol y(t_\star)$ 是该时刻状态。温度到达阈值可取 $g=T-T_\star$，单位 K；落地可取高度 $g=z$，单位 m。事件函数最好连续；求解器通常通过步内符号变化检测事件，单个大步内多次过零可能漏检。因此快速振荡、碰撞等问题要适当限制最大步长，并在不连续处停止、更新状态后重新启动积分。

\begin{equation}
\delta_Q=\frac{\max_t |Q(\boldsymbol y(t))-Q(\boldsymbol y_0)|}{Q_{\mathrm{scale}}}.
\end{equation}
$Q$ 是理论上守恒的量，如机械能 J、总质量 kg；$Q_{\mathrm{scale}}>0$ 是同单位的参考尺度；$\delta_Q$ 是无量纲守恒偏差。只有模型确实守恒时才能这样检查。例：无阻尼自由振动用总机械能检查；有阻尼时能量本来应下降，不能误判为程序错误。真实系统还有持续输入时，应比较累计输入与状态增量。

\subsection{可运行冷却模板与逐行解析}
下面程序无需数据文件：它同时验证温度解析解和阈值事件时间。所有非空代码行右侧均附中文解释；数组从下标 0 开始，所以 \texttt{y[0]} 是第一个状态。\texttt{def} 创建函数，\texttt{return} 把结果交给求解器，\texttt{assert} 是验证条件不满足就停止的检查。\texttt{t\_eval} 只是要求输出在哪些时刻，不是把内部积分步长固定下来。事件提前终止后，最后的输出采样点可能早于真实事件；事件时间要读 \texttt{t\_events}。

完整代码与每行中文注释见末尾“Python 入门与完整逐行注释模板”章的 ODE 小节。

\subsection{从模板换到自己的赛题}
先把 \texttt{rhs} 中的右端改成你的物理变化率，再按相同顺序给初始状态。两状态系统必须返回两个变化率，顺序始终对应。把参数值和单位写在一起，避免把分钟混成秒。如果环境温度随时间分段变化，可在右端计算环境温度函数；已知跳变时刻最好分段积分。对非刚性问题先用 RK45，对疑似刚性问题比较 Radau 或 BDF。最后报告算法、相对/绝对容限、最大步长、事件规则及验证结果，而不是只给一张曲线。

\chapter{一维热传导有限差分：把一根杆切成小格子}
\section{模型从守恒定律开始}
\subsection{傅里叶导热定律与方向}
\begin{equation}
q_x''=-k\frac{\partial T}{\partial x}.
\end{equation}
$q_x''$ 是沿正 $x$ 方向的热流密度，单位 W/m$^2$；双撇号只是热流密度惯用记号，不是二阶导数；$k>0$ 是导热系数 W/(m$\cdot$K)；$T(x,t)$ 是温度 K 或摄氏度；$x$ 是位置 m；$\partial$ 表示只改变一个自变量的偏导数。负号说明热从高温流向低温。例：沿右方每米降温 10 K，且 $k=2$，则向右热流密度为 20 W/m$^2$。

\subsection{小体积的能量账本}
\begin{equation}
\rho c_p A\Delta x\frac{\partial T}{\partial t}
\approx A q_x''(x,t)-A q_x''(x+\Delta x,t)+\dot q_v A\Delta x.
\end{equation}
$\rho$ 是密度 kg/m$^3$；$c_p$ 是比热容 J/(kg$\cdot$K)；$A$ 是恒定横截面积 m$^2$；$\Delta x$ 是小段长度 m；$\dot q_v$ 是单位体积发热率 W/m$^3$，上点属于记号；$t$ 是时间 s；其他量同前。左侧是储热功率，右侧是流入减流出加内部发热，三项单位都是 W。这里假定截面内温度均匀、侧面不漏热、不发生相变和显著形变；若有侧面对流，要另加损失项。

\begin{equation}
\rho c_p\frac{\partial T}{\partial t}=\frac{\partial}{\partial x}\left(k\frac{\partial T}{\partial x}\right)+\dot q_v,
\qquad \alpha=\frac{k}{\rho c_p}.
\end{equation}
$\alpha$ 是热扩散率，单位 m$^2$/s；其余符号同前。先对小段能量式除以体积，再令段长趋于零，就得到这条偏微分方程。导热系数随位置或温度变化时，$k$ 必须留在空间导数里面。例：两个材料接触时，不能拿一个平均 $k$ 随意替换整个空间，界面的热流连续性必须保持。

\begin{equation}
\frac{\partial T}{\partial t}=\alpha\frac{\partial^2 T}{\partial x^2}+s(x,t),\qquad s=\frac{\dot q_v}{\rho c_p}.
\end{equation}
$s$ 是温度源项，单位 K/s；$\partial^2 T/\partial x^2$ 是温度空间曲率，单位 K/m$^2$；$\alpha$ 与前同。此形式假设 $k,\rho,c_p$ 为常数。曲率是正时该点比两旁低，导热使它升温；曲率为负时使它降温。例：杆中央形成尖锐热峰，中央会逐渐冷却，热量向两侧扩散。

\subsection{初始条件与三类边界条件}
\begin{align}
T(x,0)&=T_{\mathrm{init}}(x),\quad 0\leq x\leq L,\\
T(0,t)&=T_L(t),\qquad T(L,t)=T_R(t).
\end{align}
$L$ 是杆长 m；$T_{\mathrm{init}}$ 是初始温度分布；$T_L,T_R$ 是左右端温度，单位同 $T$。第一行是初始条件，第二行是第一类（Dirichlet）边界。必须同时知道初始和边界信息，才能确定演化。例：两端恒为 0 摄氏度就令 $T_L=T_R=0$。为了得到光滑解，初温端点应与开始时的边界温度一致，否则会有起始边界层，可能降低观察到的收敛阶。

\begin{equation}
-k\frac{\partial T}{\partial n}=q_n'',\qquad
\frac{\partial T}{\partial n}=0\ \text{（绝热）}.
\end{equation}
$n$ 表示边界向外的法向方向，不是时间编号；$q_n''$ 是向外流出的热流密度 W/m$^2$；$\partial T/\partial n$ 是沿外法向的温度梯度 K/m；$k$ 同前。这是第二类（Neumann）边界。杆左端外法向指向负 $x$，右端指向正 $x$，因此左右端的符号不能机械照抄。绝热不是温度为零，而是没有热流穿过边界。

\begin{equation}
-k\frac{\partial T}{\partial n}=h_c(T_s-T_\infty).
\end{equation}
$h_c$ 是表面对流换热系数 W/(m$^2\cdot$K)；$T_s$ 是端面温度；$T_\infty$ 是流体环境温度；其余符号同前。这是第三类（Robin）边界。$h_c$ 不是时间步长。例：端面比环境高 20 K，$h_c=10$，则向外散热 200 W/m$^2$。辐射显著时还需加入辐射热流，方程会因四次方而非线性。

\section{空间离散与显式格式}
\begin{equation}
x_i=i\Delta x,\quad \Delta x=L/M,\quad t_n=n\Delta t,\quad T_i^n\approx T(x_i,t_n).
\end{equation}
$M$ 是空间区间数，无量纲，节点共有 $M+1$ 个；$i=0,\ldots,M$ 是空间编号；$n$ 是时间编号；$\Delta x$ 的单位 m；$\Delta t$ 的单位 s；$T_i^n$ 是温度近似值，上标 $n$ 是时间层而不是乘方。只有 $i=1,\ldots,M-1$ 是内部节点；两端由边界条件决定。例：1 m 长杆切成 10 段，需要存 11 个温度。

\begin{equation}
\left.\frac{\partial T}{\partial t}\right|_i^n\approx\frac{T_i^{n+1}-T_i^n}{\Delta t},\qquad
\left.\frac{\partial^2T}{\partial x^2}\right|_i^n\approx\frac{T_{i-1}^n-2T_i^n+T_{i+1}^n}{\Delta x^2}.
\end{equation}
所有符号同上一式；竖线表示在这个网格点计算。第一个差商单位 K/s，第二个单位 K/m$^2$。时间前向差分一阶精度，空间中心二阶差分在温度足够光滑时二阶精度。例：一排温度为 0、10、0，中央的差分曲率为负，因此中央向外散热。

\begin{equation}
r=\frac{\alpha\Delta t}{\Delta x^2},\qquad
T_i^{n+1}=rT_{i-1}^n+(1-2r)T_i^n+rT_{i+1}^n+\Delta t\,s_i^n.
\end{equation}
$r$ 是无量纲扩散比；$s_i^n=s(x_i,t_n)$ 是温度源项 K/s；其他符号同前。这就是显式 FTCS 格式。新温度由旧时刻三个温度算出，不能边计算边覆盖旧数组，否则会混用时间层。例：$r=0.4$、无热源、旧温度为 0、10、0，则新中央温度为 2。若本应平滑扩散却出现越来越大的正负交替，优先检查步长稳定性。

\begin{equation}
0\leq r\leq\frac12\quad\Longleftrightarrow\quad
\Delta t\leq\frac{\Delta x^2}{2\alpha}.
\end{equation}
这是常系数一维纯扩散、均匀网格、标准边界处理下显式格式的稳定性条件；符号与前相同。满足条件时无热源更新权重非负且和为 1，可看成三个旧温度的加权平均。空间步长缩小到一半，允许的时间步长会缩小到四分之一。增加对流、反应或特殊边界后，要重新检查相应约束，不能直接套用这一条作为全部稳定条件。

\begin{equation}
G(\xi)=1-4r\sin^2(\xi/2),\qquad |G(\xi)|\leq1.
\end{equation}
$G$ 是单步误差放大因子，无量纲；$\xi$ 是网格上的无量纲波数（弧度视为无量纲）；$\sin$ 是正弦函数。它来自把误差分解成正弦波分析。所有波数都不放大时才稳定；最坏情况 $\sin^2(\xi/2)=1$ 就得到 $r\leq1/2$。这解释了为什么某些棋盘状高频误差最先爆炸。

\section{隐式、Crank--Nicolson 与矩阵求解}
\subsection{后向欧拉：新时间层的三个点一起解}
\begin{equation}
-rT_{i-1}^{n+1}+(1+2r)T_i^{n+1}-rT_{i+1}^{n+1}
=T_i^n+\Delta t\,s_i^{n+1}.
\end{equation}
$T_i^n,r,\Delta t,s_i^{n+1}$ 均按此前定义。右边已知，左边含相邻新温度，因此所有内部点要组成联立方程一起求。常系数线性扩散下它无条件稳定，但时间精度仍是一阶；大步长会把真实变化抹得过平。稳定仅表示不爆炸，不表示答案准确。

\subsection{Crank--Nicolson：平均新旧两层的导热}
\begin{align}
-\frac r2T_{i-1}^{n+1}+(1+r)T_i^{n+1}-\frac r2T_{i+1}^{n+1}
&=\frac r2T_{i-1}^n+(1-r)T_i^n+\frac r2T_{i+1}^n\\
&\quad+\frac{\Delta t}{2}(s_i^n+s_i^{n+1}).
\end{align}
符号同前。对于光滑解、光滑源项和一致的二阶边界处理，时间与空间均可达到二阶精度。它对线性扩散无条件稳定，但大时间步可能产生不衰减得足够快的交替振荡，尤其是阶跃初温。遇到这种现象，应减小时间步，或先用若干隐式欧拉小步平滑起始层，而不是仅仅宣布“无条件稳定所以没问题”。

\subsection{一个参数统一三种方法}
\begin{align}
-\theta rT_{i-1}^{n+1}+(1+2\theta r)T_i^{n+1}-\theta rT_{i+1}^{n+1}
&=T_i^n+(1-\theta)r(T_{i-1}^n-2T_i^n+T_{i+1}^n)\\
&\quad+\Delta t[(1-\theta)s_i^n+\theta s_i^{n+1}].
\end{align}
$\theta\in[0,1]$ 是无量纲时间加权系数；$\theta=0$ 为显式，$\theta=1$ 为隐式欧拉，$\theta=1/2$ 为 Crank--Nicolson；其余符号同前。代码用这个统一格式避免重复三份更新逻辑。下面程序虽然显式情形也调用带状线性求解器，但这时矩阵是单位阵，数学结果与直接更新相同。

\begin{equation}
B\boldsymbol u^{n+1}=\boldsymbol b^n,\qquad
B=\operatorname{tridiag}(-\theta r,1+2\theta r,-\theta r).
\end{equation}
$\boldsymbol u^{n+1}=(T_1^{n+1},\ldots,T_{M-1}^{n+1})^{\mathsf T}$ 是内部温度向量；$B$ 是无量纲的三对角系数矩阵；$\boldsymbol b^n$ 是旧温度、源项和边界项组成的右端，单位 K；$\operatorname{tridiag}$ 表示按下对角、主对角、上对角构造矩阵。不要显式计算 $B^{-1}$，应调用线性方程求解器。网格很大时利用三对角结构，比构造完整密集矩阵节省大量存储。

\subsection{非零端点一定要补到右端}
\begin{align}
b_1^n&\mathrel{+}=\theta rT_L(t_{n+1}),\\
b_{M-1}^n&\mathrel{+}=\theta rT_R(t_{n+1}).
\end{align}
$b_1^n$ 和 $b_{M-1}^n$ 是右端向量首尾元素，单位 K；符号 $\mathrel{+}=$ 表示“在已有值上再加”，不是新的物理定律；$T_L,T_R$ 为给定边界温度。旧时刻边界贡献已经包含在统一格式的旧时刻差分项中，新时刻边界贡献还必须由左侧移到右侧。下面代码的两端均为零，故这两项恰好为零；直接改成非零边界却不加它们会得到错误结果。

\subsection{绝热与对流边界的离散例子}
\begin{equation}
\frac{-3T_0+4T_1-T_2}{2\Delta x}=0\quad\text{（左端绝热）}.
\end{equation}
$T_0,T_1,T_2$ 是同一时间层从左端开始的三个温度；$\Delta x$ 同前。这是左端一阶导数的二阶单边差分，导数为零即绝热。它要作为边界方程并入线性系统，不能既把 $T_0$ 固定为常数又声称绝热。更简单的 $T_0=T_1$ 只有一阶边界近似，可能影响总体收敛阶。

\begin{equation}
-k\frac{3T_M-4T_{M-1}+T_{M-2}}{2\Delta x}
=h_c(T_M-T_\infty)\quad\text{（右端对流）}.
\end{equation}
$T_M$ 是右端温度，$T_{M-1},T_{M-2}$ 是其左侧两个点；$k,h_c,T_\infty$ 同前。这是右端外法向为正 $x$ 的二阶单边形式。左端不能直接复制负号，因为外法向反向。边界处理改变了系数矩阵结构，必须相应修改求解器存储；本页是建模公式，后面的基础模板只实现两端恒温。

\section{材料变化、侧面散热与验证}
\begin{equation}
\rho_i c_{p,i}\frac{T_i^{n+1}-T_i^n}{\Delta t}
=\frac{k_{i+1/2}(T_{i+1}^n-T_i^n)-k_{i-1/2}(T_i^n-T_{i-1}^n)}{\Delta x^2}+\dot q_{v,i}^n.
\end{equation}
$\rho_i,c_{p,i}$ 是第 $i$ 点的材料参数；$k_{i+1/2}$ 是 $i$ 与 $i+1$ 两点之间界面的有效导热系数；半整数标记界面，不是半个数组元素；其余符号与前同。左右都是 W/m$^3$。这是通量守恒型离散，适合空间材料参数变化；该式是显式时间处理，仍需专门确定稳定步长。界面有接触热阻时要另加热阻，不能假装完美接触。

\begin{equation}
k_{i+1/2}=\frac{2k_i k_{i+1}}{k_i+k_{i+1}}.
\end{equation}
$k_i,k_{i+1}$ 是相邻两材料的导热系数，单位 W/(m$\cdot$K)；$k_{i+1/2}$ 是同单位界面系数。等长半单元、分段常系数、完美接触情况下，串联热阻给出调和平均。例：$k_i=1$、$k_{i+1}=100$，有效系数约为 1.98，而不是算术平均 50.5，因为差导体限制了整个通路。

\begin{equation}
\rho c_p A\frac{\partial T}{\partial t}
=kA\frac{\partial^2T}{\partial x^2}-h_cP(T-T_\infty)+A\dot q_v.
\end{equation}
$P$ 是截面周长 m；$A$ 是截面积 m$^2$；其余符号同前。每项单位 W/m，这是细长杆侧面对流散热的一维模型。适用于每个截面内温差较小的情况。侧面散热是沿整根杆持续损失，与只有端面散热不同；若已经把侧面项写进 PDE，不要在同一物理表面重复计入。

\subsection{正弦解析解：三个方法共用的检验题}
\begin{equation}
T(x,t)=T_b+A_T\sin\left(\frac{\pi x}{L}\right)
\exp\left[-\alpha\left(\frac{\pi}{L}\right)^2t\right].
\end{equation}
$T_b$ 是两端共同恒温；$A_T$ 是初始波峰超过边界温度的幅值 K；$\pi$ 是圆周率；$L,x,\alpha,t$ 如前。适用条件是常系数、无热源、初温恰为该正弦形、两端恒为 $T_b$。指数无量纲。代码取 $T_b=0$ 摄氏度、$A_T=10$ K、$L=1$ m、$\alpha=0.01$ m$^2$/s。它是人为构造的教学验证题，不是某种具体金属的实验测量。

\begin{equation}
E_\infty=\max_i|T_i^{\mathrm{num}}-T_i^{\mathrm{exact}}|,\qquad
E_2=\left(\frac{1}{M+1}\sum_{i=0}^M(T_i^{\mathrm{num}}-T_i^{\mathrm{exact}})^2\right)^{1/2}.
\end{equation}
$T_i^{\mathrm{num}}$ 是某一固定时刻的数值温度，$T_i^{\mathrm{exact}}$ 是同一坐标、同一时刻的解析温度；$E_\infty$ 是最大误差，$E_2$ 是节点均方根误差，单位均为 K；$M$ 是区间数。这里 $E_2$ 是离散节点平均，不能与没有归一化的连续积分范数混为一谈。

\begin{equation}
E(\Delta x,\Delta t)\approx C_x\Delta x^2+C_t\Delta t^q,
\quad q=1\ \text{（欧拉）},\quad q=2\ \text{（CN）}.
\end{equation}
$E$ 是温度误差；$C_x$ 的单位 K/m$^2$，$C_t$ 的单位 K/s$^q$；$q$ 是时间阶；其他符号同前。此为光滑问题的主导误差尺度描述，实际有符号误差可能抵消。代码让 $\Delta t$ 与 $\Delta x^2$ 同阶，所以三个方法的联合加密误差都约按 $\Delta x^2$ 缩小；不能据此声称隐式欧拉本身时间二阶。单独检验时间阶，应固定非常细的空间网格，只减时间步；单独检验空间阶则应把时间误差压得足够小。

\begin{equation}
\mathcal E(t)=\rho c_pA\int_0^L[T(x,t)-T_{\mathrm{ref}}]\,\mathrm dx.
\end{equation}
$\mathcal E$ 是相对于参考温度 $T_{\mathrm{ref}}$ 的储热量 J；$T_{\mathrm{ref}}$ 是固定温度，单位同 $T$；其余符号同前。两端绝热、无热源时它应守恒；两端恒温时它通常不守恒，因为热能从端点流出。验证时要先写清楚系统的能量边界，再判断能量变化对不对。

\subsection{可运行三方法模板与逐行解析}
程序对显式、隐式和 CN 各计算两套网格，与同一个解析解比较，检查误差随加密降低、端点不漂移、无热源时温度不超过初始最大值。每行的中文注释解释了数据和运算。\texttt{u[1:-1]} 表示排除两个端点的内部温度；\texttt{u[:-2]}、\texttt{u[2:]} 分别是这些内部点的左邻、右邻；\texttt{**2} 表示平方。\texttt{solve\_banded((1,1),...)} 中两个 1 表示只有一条下对角线和一条上对角线。此模板固定零温度边界、无源项、常材料参数；用于其他条件时，要按本章对应方程修改并重新验证。

完整代码与每行中文注释见末尾“Python 入门与完整逐行注释模板”章的热传导小节。

\section{本部分的来源与复现说明}
经典差分教材的出版方入口：\url{https://epubs.siam.org/doi/10.1137/1.9780898717938}；经典 ODE 教材出版方入口：\url{https://link.springer.com/book/10.1007/978-3-540-78862-1}。
上述方程推导和示例为教学性整理；软件行为核对于 2026 年 9 月 10 日。SciPy 官方文档支持 \texttt{solve\_ivp} 的方法选择、误差容限和事件接口：\url{https://docs.scipy.org/doc/scipy/reference/generated/scipy.integrate.solve_ivp.html}。带状线性方程求解器的存储格式见：\url{https://docs.scipy.org/doc/scipy/reference/generated/scipy.linalg.solve_banded.html}。软件接口可能随版本变化，提交比赛代码时应记录 Python、NumPy 和 SciPy 版本。热方程基本差分格式可参照经典教材：John C. Strikwerda, \textit{Finite Difference Schemes and Partial Differential Equations}, 2nd ed., SIAM, 2004；ODE 方法可参照 Hairer, N\o rsett and Wanner, \textit{Solving Ordinary Differential Equations I}, 2nd revised ed., Springer, 1993。书目用于进一步学习，不意味着本手册覆盖了这些书的全部定理和证明。

\chapter{基础定理：为什么方程能解、算法可信}
本章是条件清单式的入门解释，不是完整证明。使用定理前，必须核查它的假设；没有核查就引用定理，不能给模型增加可信度。
\section{Picard--Lindelöf：初值能否决定唯一轨迹}
\begin{equation}\dot{\boldsymbol y}=\boldsymbol f(t,\boldsymbol y),\quad\boldsymbol y(t_0)=\boldsymbol y_0,\quad \|\boldsymbol f(t,\boldsymbol u)-\boldsymbol f(t,\boldsymbol v)\|\le L\|\boldsymbol u-\boldsymbol v\|.\end{equation}
$t$ 为时间，$t_0$ 为初始时刻，$\boldsymbol y$ 为状态，$\boldsymbol y_0$ 为初始状态；$\boldsymbol f$ 是变化率；$\boldsymbol u,\boldsymbol v$ 是初始状态附近的两个候选状态；$\|\cdot\|$ 是一致选用的向量范数；$L\ge0$ 为局部 Lipschitz 常数。如果状态用同一单位且时间为 s，$L$ 用 s$^{-1}$。异单位状态应先缩放。

当 $\boldsymbol f$ 在初值附近连续，并对状态满足局部 Lipschitz 条件（在局部时间区间上一致）时，存在初始时刻附近的唯一解。白话是“变化规律不能在相邻状态间突然变得过于陡峭”。它只保证局部，不保证解永远存在；也不等于方程一定有初等函数解析解。基础来源：\url{https://epubs.siam.org/doi/10.1137/1.9780898719222}。

\section{Gronwall 型估计：小误差可能随时间放大}
\begin{equation}e(t)\le e_0+L\int_0^t e(s)\,\mathrm ds\quad\Longrightarrow\quad e(t)\le e_0e^{Lt},\quad t\ge0.\end{equation}
$e(t)\ge0$ 是连续误差大小，$e_0\ge0$ 为初始误差上界，$L\ge0$ 为常数，$s$ 为积分用的时间变量；右端指数中的 $e$ 是自然常数，和误差函数 $e(t)$ 的字母重用但含义不同。若 $t$ 用 s，则 $L$ 用 s$^{-1}$。误差上界与误差同单位。白话是“即使一开始差得很少，长期结果也可能不再接近”；上界可能很保守，不代表实际误差一定按指数长大。它常用于 ODE 稳定性和收敛分析，可结合前节 ODE 教材学习。

\section{散度定理：体内守恒怎样变成边界通量}
\begin{equation}\int_\Omega\nabla\cdot\boldsymbol F\,\mathrm dV=\int_{\partial\Omega}\boldsymbol F\cdot\boldsymbol n\,\mathrm dA.\end{equation}
$\Omega$ 是三维体积区域，$\partial\Omega$ 是其封闭边界，$\boldsymbol F$ 是足够光滑的向量通量场，$\boldsymbol n$ 是向外单位法向量，$\mathrm dV$ 和 $\mathrm dA$ 分别为微小体积和面积；点号为点积。若 $\boldsymbol F$ 是热通量 W/m$^2$，两边都是 W。区域边界通常要求分片光滑。

白话是“内部各小块的净流出加起来，只剩整个外边界的流出”；内部公共面的相反通量会抵消。它连接热传导、流体守恒与有限体积法。法向量选成向内时符号会反过来。来源：\url{https://ocw.mit.edu/ans7870/18/18.013a/textbook/HTML/chapter22/section03.html}。

\section{分部积分：把导数从嘈杂数据上移开}
\begin{equation}\int_a^b u'(x)v(x)\,\mathrm dx=[u(x)v(x)]_a^b-\int_a^b u(x)v'(x)\,\mathrm dx.\end{equation}
$u,v$ 是足够光滑的函数，$u',v'$ 是对 $x$ 的导数，$a,b$ 是积分端点，方括号表示上端值减下端值。若 $u$ 单位为温度、$v$ 无量纲、$x$ 用米，则各项都是温度。选择两端为零的 $v$ 可消掉边界项。有限元和弱形式模型识别会用它；边界项只能因实际条件为零而去掉，不能随意丢弃。

\section{Lax 等价定理：一致性与稳定性怎样合起来}
\begin{equation}\text{一致的线性差分格式：}\qquad \text{稳定}\ \Longleftrightarrow\ \text{收敛}.\end{equation}
这行没有物理参数，是逻辑关系；双向箭头表示在定理条件下两性质等价。一致性是步长趋零时差分逼近原方程；稳定性是有限时间内误差放大有与网格无关的界；收敛是网格趋细时数值解趋近真实解。经典表述针对适定线性初值问题及相应一致的线性差分近似，使用相容的函数空间和范数。

不要把这句简写推广到任意非线性算法，也不要把“某次没爆炸”当成稳定性证明。热方程显式法需要同时满足稳定步长与一致离散。基础课程入口：\url{https://www-math.mit.edu/~stevenj/18.336/}。

\section{热方程最大值原理：没有热源就不会凭空更热}
\begin{equation}\min_{\Gamma_p}T\le T(x,t)\le\max_{\Gamma_p}T,\quad \Gamma_p=\{(x,0):0\le x\le L\}\cup\{(0,t),(L,t):0\le t\le t_f\}.\end{equation}
$T$ 是满足一维无源热方程的连续温度，$x$ 为杆内位置，$t$ 为时间，$L>0$ 是长度，$t_f>0$ 是终止时间；$\Gamma_p$ 是“初始时刻整根杆加上之后两端”的抛物边界，$\cup$ 表示集合并集。对于常正热扩散率、适当光滑解和 Dirichlet 边界，内部温度不会超过这组初始与边界数据的最大值，也不会低于最小值。所有温度单位一致。

若两端恒温且初温处于某范围，数值温度大幅超出该范围通常应检查实现或离散振荡。有热源或不同方程时需修改结论。Crank--Nicolson 的稳定性不自动保证所有步长都保持离散最大值性质。参见差分章节的 Strikwerda 教材。

\chapter{非线性优化：把“最好”写成可计算的问题}
\section{目标、约束和决策变量}
\subsection{约束优化标准式}
\begin{equation}\min_{\boldsymbol x\in\mathbb R^n}f(\boldsymbol x),\quad g_i(\boldsymbol x)\le0\ (i=1,\ldots,m),\quad h_j(\boldsymbol x)=0\ (j=1,\ldots,p).\end{equation}
$\boldsymbol x$ 是允许调节的 $n$ 维决策向量，如厚度、功率、时间；$f$ 是要最小化的单个目标，如成本；$g_i$ 是 $m$ 条不等式约束，$h_j$ 是 $p$ 条等式约束；$\mathbb R^n$ 表示 $n$ 个实数组成的空间。每条约束内部必须量纲一致，各变量可以有各自单位。最大化收益可以最小化其负数。参数是给定或估计值，决策变量才是优化器主动改变的量。
\subsection{上下界}
\begin{equation}l_k\le x_k\le u_k,\quad k=1,\ldots,n.\end{equation}
$x_k$ 为第 $k$ 个变量，$l_k,u_k$ 为其物理下上界，单位与 $x_k$ 相同；$n$ 是变量数。厚度非负不意味着允许零厚度，真实工程可要求严格正的最小加工厚度。
\subsection{目标缩放}
\begin{equation}\widetilde f(\boldsymbol x)=\frac{f(\boldsymbol x)}{f_0},\qquad \widetilde x_k=\frac{x_k-l_k}{u_k-l_k}.\end{equation}
$f_0>0$ 是目标典型尺度，$u_k>l_k$，波浪号表示无量纲缩放。这样能减少变量相差多个数量级导致的数值困难。结果必须换回原单位。固定变量上下界相同时，不用该缩放式，应从未知量中移除。

\section{最基本的下降方法}
本节两式针对已无量纲化的变量和目标；梯度、Hessian、方向及步长均按无量纲量理解。
\subsection{梯度下降}
\begin{equation}\boldsymbol x_{k+1}=\boldsymbol x_k-\alpha_k\nabla f(\boldsymbol x_k).\end{equation}
$k$ 为迭代编号，$\boldsymbol x_k$ 为当前候选解，$\nabla f$ 为偏导组成的梯度，$\alpha_k>0$ 是步长。沿负梯度通常在足够小步长下下降；本式是无约束迭代，有约束不能直接照搬。建议先无量纲化；步长过大可能发散。
\subsection{牛顿方向}
\begin{equation}H_k\boldsymbol d_k=-\nabla f(\boldsymbol x_k),\qquad \boldsymbol x_{k+1}=\boldsymbol x_k+\alpha_k\boldsymbol d_k.\end{equation}
$H_k=\nabla^2 f(\boldsymbol x_k)$ 是 Hessian，即二阶偏导矩阵，$\boldsymbol d_k$ 为方向，$\alpha_k$ 为步长。用线性求解获得方向；$H_k$ 非正定时方向未必下降，实际算法需修正 Hessian 或用信赖域。不要声称任意非线性问题都能由牛顿法求得全局最优。

\section{KKT 条件与凸性}
\subsection{拉格朗日函数}
\begin{equation}\mathcal L(\boldsymbol x,\boldsymbol\lambda,\boldsymbol\nu)=f(\boldsymbol x)+\sum_{i=1}^m\lambda_i g_i(\boldsymbol x)+\sum_{j=1}^p\nu_j h_j(\boldsymbol x).\end{equation}
$\mathcal L$ 是把目标与约束合起来的函数；$\lambda_i$ 为不等式乘子，$\nu_j$ 为等式乘子；$f,g_i,h_j,m,p,\boldsymbol x$ 含义同标准式。乘子的单位应使每项与 $f$ 同单位。乘子可辅助解释资源限制的边际价值，但符号解释取决于约束写法。
\subsection{Karush--Kuhn--Tucker 条件}
\begin{equation}\nabla_x\mathcal L=0,\quad g_i\le0,\quad h_j=0,\quad\lambda_i\ge0,\quad\lambda_i g_i=0.\end{equation}
所有函数与梯度都在候选最优点计算；$i=1,\ldots,m$，$j=1,\ldots,p$。依次是驻点、原始可行性、对偶可行性和互补松弛。最后一项表示“不起限制作用的约束，其乘子必须为零”。对于可微问题，局部最优点在满足适当约束资格条件（例如 LICQ）时满足 KKT；任意问题的 KKT 点都全局最优是错误说法。LICQ 指所有等式约束和取等号的不等式约束的梯度向量线性无关，通俗说这些局部限制不能互相重复。等式乘子可以为正、为负或为零，不要求非负。
\subsection{凸函数定义}
\begin{equation}f(\tau\boldsymbol x+(1-\tau)\boldsymbol y)\le\tau f(\boldsymbol x)+(1-\tau)f(\boldsymbol y),\qquad 0\le\tau\le1.\end{equation}
$\boldsymbol x,\boldsymbol y$ 是凸定义域内任意两点，$\tau$ 是无量纲混合权重。意思是连接两点的弦在函数图像上方。若目标与不等式函数凸、等式仿射，则是凸优化，局部最优也是全局最优；满足 KKT 的点为全局最优。Slater 条件等可保障凸问题强对偶和 KKT 必要性。严格凸目标在凸可行域上至多一个最优点，但并不保证最优点存在。参考 Boyd 与 Vandenberghe 的教材网站：\url{https://web.stanford.edu/~boyd/cvxbook/}。

\section{参数拟合：和设计优化有什么不同}
\subsection{加权非线性最小二乘}
\begin{equation}\widehat{\boldsymbol\theta}=\arg\min_{\boldsymbol\theta}\sum_{i=1}^N\left[\frac{f(t_i;\boldsymbol\theta)-y_i}{s_i}\right]^2.\end{equation}
$t_i$ 是第 $i$ 次观测的时间或输入，$y_i$ 为测量结果，$f$ 为预测函数，$\boldsymbol\theta$ 是待估参数，$s_i>0$ 为该观测误差标准差或明确指定的权重尺度，$N$ 是样本数；帽子表示估计结果。残差除以 $s_i$ 后无量纲。观测误差独立且正态、标准差已知时，它对应最大似然估计；相关误差需要协方差加权。
\subsection{Tikhonov 正则化}
\begin{equation}\min_{\boldsymbol\theta}\|\boldsymbol r(\boldsymbol\theta)\|_2^2+\lambda\|L(\boldsymbol\theta-\boldsymbol\theta_0)\|_2^2.\end{equation}
$\boldsymbol r$ 是残差向量，$\lambda\ge0$ 为正则化强度，$L$ 是选定的缩放或差分矩阵，$\boldsymbol\theta_0$ 是参考参数，$\|\cdot\|_2$ 是欧氏范数。建议两项先无量纲化。第二项抑制过大或不平滑参数，不能凭空增加观测信息；$\lambda$ 应通过验证或误差原则选取。
\subsection{Gauss--Newton 近似}
\begin{equation}(J^\mathsf TJ)\Delta\boldsymbol\theta=-J^\mathsf T\boldsymbol r,\qquad J_{ij}=\frac{\partial r_i}{\partial\theta_j}.\end{equation}
$J$ 是残差对参数的 Jacobian，$i$ 编号残差、$j$ 编号参数，$\Delta\boldsymbol\theta$ 为参数修正。该式忽略残差的二阶导数贡献，适合残差较小等情况。数值实现宜用 QR、SVD 或成熟最小二乘求解器；直接构造正规方程会恶化条件数。数据不足或参数作用相似时，即使拟合曲线很好，参数也可能不唯一。

\section{多目标、鲁棒和概率约束}
\subsection{多目标加权和}
\begin{equation}\min_{\boldsymbol x}\sum_{j=1}^q w_j\frac{f_j(\boldsymbol x)-a_j}{s_j},\qquad w_j\ge0,\quad\sum_{j=1}^q w_j=1.\end{equation}
$q$ 是目标数，$f_j$ 是第 $j$ 个目标，$a_j$ 是参考值，$s_j>0$ 是同单位尺度，$w_j$ 是无量纲偏好权重。先把所有目标统一为越小越好。权重体现偏好而非客观真理；非凸 Pareto 前沿的部分点无法通过加权和取得。
\subsection{最坏情形鲁棒优化}
\begin{equation}\min_{\boldsymbol x}\max_{\boldsymbol\theta\in\mathcal U} f(\boldsymbol x,\boldsymbol\theta).\end{equation}
$\boldsymbol x$ 是决策，$\boldsymbol\theta$ 是不确定参数，$\mathcal U$ 是有依据设定的不确定集合，$f$ 是代价。意思是在可能的参数情形中也尽量控制最坏损失。集合过大可能过分保守，必须解释范围来源；约束也可要求对整个集合成立。
\subsection{机会约束}
\begin{equation}\mathbb P\{g(\boldsymbol x,\boldsymbol\xi)\le0\}\ge1-\varepsilon.\end{equation}
$\boldsymbol\xi$ 是具有指定概率分布的随机输入，$g$ 是约束函数，$\varepsilon\in(0,1)$ 是允许违规概率，$\mathbb P$ 为概率。比如温度超限概率不超过 1\%，则 $\varepsilon=0.01$。有限样本中的合格比例不是自动的总体概率保证。

\section{一个可以手算核验的优化例子}
\begin{equation}\min_{x,y}(x-2)^2+(y-1)^2,\quad x\ge0,\ y\ge0,\ x+y\le2.\end{equation}
$x,y$ 是无量纲决策量。无约束最优 $(2,1)$ 不可行；约束边界上令 $y=2-x$，最优是 $(1.5,0.5)$，目标值 0.5。几何意思是找可行三角形中离点 $(2,1)$ 最近的点。这是用于验证代码的凸问题，不能代表真实非凸问题的难度。

\section{模板怎么选、结果怎么检查}
用 \texttt{scipy.optimize.least\_squares} 做残差拟合；用 \texttt{minimize} 的 SLSQP 处理平滑目标和等式、不等式约束。SciPy 旧式字典约束的 \texttt{ineq} 表示函数大于等于零，与本章标准式的小于等于零方向相反。因此 $x+y\le2$ 要写为 \texttt{2-x-y}。接口依据：\url{https://docs.scipy.org/doc/scipy/reference/generated/scipy.optimize.minimize.html}。

验收顺序：读求解器成功标志和信息；重新计算目标；检查每个约束残差与上下界；改初值重复；改变缩放与容差；检查小扰动下结论是否稳定。多初值只能提高找到较好解的机会，不能证明全局最优。附录中的优化模板还演示了非线性指数拟合；其中无噪声合成数据只用于教学。

\chapter{物理建模工具箱：先判断现象，再选公式}
本章是面向国赛工程与物理建模题的通用工具箱，不代表任何一届题目的押题清单。一个物理公式只有在假设成立时才有用。读题后先画出研究对象和边界，再标出质量、力、热量或流量从哪里进入、从哪里离开。所有算例数据均为教学假设值，不是比赛数据或材料推荐值。除特别说明外，使用国际单位制；角度代入程序前换成弧度，温差可用摄氏度差，绝对温度用开尔文。

\section{力学：物体为什么运动、变形、振动}
\subsection{M1：位置、速度、加速度}
\begin{equation}
\mathbf v=\frac{d\mathbf r}{dt},\qquad \mathbf a=\frac{d\mathbf v}{dt}=\frac{d^2\mathbf r}{dt^2}.
\end{equation}
\textbf{符号：} $t$ 是时间（s）；$\mathbf r$ 是相对选定原点的位置向量（m）；$\mathbf v$ 是速度向量（m/s）；$\mathbf a$ 是加速度向量（m/s$^2$）；$d/dt$ 表示对时间求导，粗体表示有大小和方向。
\textbf{白话：} 位置改变得多快叫速度，速度改变得多快叫加速度。\textbf{条件与误用：} 轨迹需要足够光滑；速率是速度的大小，匀速圆周运动的速率不变但方向变，所以加速度不为零。用实验数据求导时，噪声会被放大，通常先平滑再求导。

\subsection{M2：匀加速直线运动}
\begin{equation}
v=v_0+at,\qquad x=x_0+v_0t+\frac12at^2.
\end{equation}
\textbf{符号：} $x,x_0$ 为当前、初始位置（m）；$v,v_0$ 为当前、初始带符号速度（m/s）；$a$ 为常加速度（m/s$^2$）；$t$ 为从初始时刻起经过的时间（s）。
\textbf{白话：} 每秒增加相同速度时才能直接套这两式。\textbf{条件与误用：} 仅适用于同一直线方向上的常加速度；有随速度变化的空气阻力时应改用微分方程。\textbf{算例：} 从静止开始，$a=2\ \mathrm{m/s^2}$，经过 $3\ \mathrm{s}$，速度为 $6\ \mathrm{m/s}$，位移为 $9\ \mathrm m$。

\subsection{M3：牛顿第二定律与静力平衡}
\begin{equation}
\sum_j\mathbf F_j=m\mathbf a;\qquad \sum_j\mathbf F_j=\mathbf0,\quad \sum_j\boldsymbol\tau_j=\mathbf0\quad\text{（静力平衡）}.
\end{equation}
\textbf{符号：} $\mathbf F_j$ 为第 $j$ 个外力（N）；$j$ 为力的编号；$\sum$ 为求和；$m$ 为不变质量（kg）；$\mathbf a$ 为质心加速度（m/s$^2$）；$\boldsymbol\tau_j$ 为该力对同一参考点的力矩（N$\cdot$m）；$\mathbf0$ 为零向量。
\textbf{白话：} 合力决定平动，合力矩决定是否有转动趋势。\textbf{条件与误用：} 使用惯性参考系；静止刚体要同时满足力和力矩平衡。变质量喷射系统不能只写固定质量的 $m\mathbf a$ 而忽略喷出物的动量流。

\subsection{M4：重力与万有引力}
\begin{equation}
F_g=mg,\qquad F=G\frac{m_1m_2}{r^2}.
\end{equation}
\textbf{符号：} $F_g,F$ 为引力大小（N）；$m,m_1,m_2$ 为质量（kg）；$g$ 为当地重力加速度（常用教学近似 $9.81\ \mathrm{m/s^2}$）；$r$ 为质点间距或球对称物体中心间距（m）；$G$ 为万有引力常数，约 $6.6743\times10^{-11}\ \mathrm{m^3/(kg\,s^2)}$。
\textbf{白话：} 地面附近用 $mg$，天体尺度用随距离变化的引力。\textbf{条件与误用：} 第二式适合质点或不重叠球对称物体；不能把物体表面间距当中心间距；地面高度变化很小时才把 $g$ 当常数。

\subsection{M5：库仑摩擦}
\begin{equation}
|F_s|\leq\mu_sN,\qquad |F_k|=\mu_kN.
\end{equation}
\textbf{符号：} $F_s,F_k$ 为静摩擦、滑动摩擦的带符号分量（N），竖线表示大小；$\mu_s,\mu_k$ 为无量纲静、动摩擦系数；$N$ 为法向接触力（N）。
\textbf{白话：} 没滑动时摩擦按需要提供，超过上限才滑动。\textbf{条件与误用：} 干接触、经验系数近似；静摩擦不是始终等于最大值；斜面、加速升降和额外压紧时 $N$ 不一定等于重力。方向反对接触面的相对滑动或其趋势。核对来源：\href{https://openstax.org/books/university-physics-volume-1/pages/6-key-equations}{OpenStax：牛顿定律应用关键公式}。

\subsection{M6：功、动能与功率}
\begin{equation}
W=\int_C\mathbf F\cdot d\mathbf r,\qquad K=\frac12mv^2,\qquad W_{\rm net}=K_2-K_1,\qquad P=\mathbf F\cdot\mathbf v.
\end{equation}
\textbf{符号：} $W,W_{\rm net}$ 为某力、合力做功（J）；$C$ 为运动路径；$\mathbf F$ 为力（N）；$d\mathbf r$ 为微小位移（m）；点号表示点积；$K,K_1,K_2$ 为当前、初态、末态平动动能（J）；$m$ 为质量（kg）；$v$ 为速率（m/s）；$P$ 为瞬时机械功率（W）；$\mathbf v$ 为速度向量（m/s）。
\textbf{白话：} 只有沿位移方向的力分量做功；功率是做功快慢。\textbf{条件与误用：} 固定质量质点的动能定理；不能把所有功都当成热，也不能把摩擦损耗计入后又重复扣除。

\subsection{M7：机械能与势能}
\begin{equation}
E=K+U_g+U_s,\qquad U_g=mgh,\qquad U_s=\frac12kx^2,\qquad E_2-E_1=W_{\rm nc}.
\end{equation}
\textbf{符号：} $E,E_1,E_2$ 为机械能及初末值（J）；$K$ 为动能（J）；$U_g,U_s$ 为重力、弹簧势能（J）；$m$ 为质量（kg）；$g$ 为重力加速度（m/s$^2$）；$h$ 为相对零势能面的高度（m）；$k$ 为线性弹簧刚度（N/m）；$x$ 为相对自然长度的伸长或压缩量（m）；$W_{\rm nc}$ 为非保守力做功（J）。
\textbf{白话：} 能量在速度、高度、弹簧变形之间转换，摩擦会减少机械能。\textbf{条件与误用：} $g$ 近似不变且弹簧线性；不是任何系统的机械能都守恒；零势能面可任选但全程一致。

\subsection{M8：动量、冲量与碰撞}
\begin{equation}
\mathbf p=m\mathbf v,\qquad \mathbf p_2-\mathbf p_1=\int_{t_1}^{t_2}\mathbf F_{\rm ext}\,dt.
\end{equation}
\textbf{符号：} $\mathbf p,\mathbf p_1,\mathbf p_2$ 为动量及初末值（kg$\cdot$m/s）；$m$ 为质量（kg）；$\mathbf v$ 为速度（m/s）；$\mathbf F_{\rm ext}$ 为系统合外力（N）；$t,t_1,t_2$ 为时间和碰撞开始、结束时刻（s）；积分结果叫冲量（N$\cdot$s）。
\textbf{白话：} 外力持续作用的累计效果改变动量。\textbf{条件与误用：} 对封闭系统统计所有物体；外冲量可忽略才有总动量守恒。非弹性碰撞可守恒动量但不守恒机械能。

\subsection{M9：圆周运动}
\begin{equation}
v=\omega R,\qquad a_n=\frac{v^2}{R}=\omega^2R,\qquad F_n=ma_n.
\end{equation}
\textbf{符号：} $v$ 为切向速率（m/s）；$\omega$ 为角速度大小（rad/s）；$R$ 为轨道半径（m）；$a_n$ 为指向圆心的法向加速度（m/s$^2$）；$F_n$ 为合力的向心分量（N）；$m$ 为质量（kg）。
\textbf{白话：} 转弯需要持续把速度方向往圆心方向拉。\textbf{条件与误用：} 圆周半径固定；向心力是已有力的合效果，不能在受力图中再额外加一股“向心力”。

\subsection{M10：转动惯量与平行轴定理}
\begin{equation}
I=\int r_\perp^2\,dm,\qquad I=I_C+md^2,\qquad K_{\rm rot}=\frac12I\omega^2.
\end{equation}
\textbf{符号：} $I,I_C$ 为指定轴、平行质心轴的转动惯量（kg$\cdot$m$^2$）；$r_\perp$ 为质量微元到轴的垂直距离（m）；$dm$ 为质量微元（kg）；$m$ 为总质量（kg）；$d$ 为两平行轴间距（m）；$K_{\rm rot}$ 为绕指定固定轴的转动动能（J）；$\omega$ 为角速度（rad/s）。
\textbf{白话：} 同样质量离轴越远，越难转起来。\textbf{条件与误用：} 刚体、轴必须明确；平行轴定理的一根轴必须经过质心。来源：\href{https://openstax.org/books/university-physics-volume-1/pages/10-key-equations}{OpenStax：固定轴转动公式}。

\subsection{M11：力矩、角动量与固定轴动力学}
\begin{equation}
\boldsymbol\tau=\mathbf r\times\mathbf F,\qquad \mathbf L=\mathbf r\times\mathbf p,\qquad \frac{d\mathbf L}{dt}=\boldsymbol\tau_{\rm ext},\qquad I\ddot\theta=\tau.
\end{equation}
\textbf{符号：} $\boldsymbol\tau,\boldsymbol\tau_{\rm ext}$ 为力矩及合外力矩（N$\cdot$m）；$\mathbf r$ 为参考点指向受力点或质点的位置向量（m）；$\mathbf F$ 为力（N）；$\times$ 为叉积；$\mathbf L$ 为角动量（kg$\cdot$m$^2$/s）；$\mathbf p$ 为动量（kg$\cdot$m/s）；$t$ 为时间（s）；$I$ 为固定轴转动惯量（kg$\cdot$m$^2$）；$\theta$ 为转角（rad）；双点为时间二阶导数；$\tau$ 为轴向合力矩（N$\cdot$m）。
\textbf{白话：} 离转轴远且方向合适的力更容易转动物体。\textbf{条件与误用：} 角动量定理采用惯性系固定参考点；最后一式仅是恒定惯量固定轴形式，三维自由转动一般要用惯量张量。

\subsection{M12：弹簧、阻尼与受迫振动}
\begin{equation}
m\ddot x+c\dot x+kx=F_0\cos(\Omega t),\qquad \omega_n=\sqrt{k/m},\qquad \zeta=\frac{c}{2\sqrt{km}}.
\end{equation}
\textbf{符号：} $m$ 为质量（kg）；$x$ 为相对平衡点位移（m）；$t$ 为时间（s）；单点、双点为速度和加速度；$c$ 为黏性阻尼系数（N$\cdot$s/m）；$k$ 为刚度（N/m）；$F_0$ 为激励力幅值（N）；$\Omega$ 为外力角频率（rad/s）；$\omega_n$ 为无阻尼固有角频率（rad/s）；$\zeta$ 为无量纲阻尼比。
\textbf{白话：} 惯性、阻尼、弹簧恢复力共同对抗外力。\textbf{条件与误用：} 小位移、线性弹簧、速度正比阻尼；$\Omega$ 接近 $\omega_n$ 时要警惕共振。必须再给初始位移、初始速度才能得到唯一运动曲线。

\subsection{M13：欠阻尼自由振动的解析核验式}
\begin{equation}
x(t)=e^{-\zeta\omega_nt}\left(A\cos\omega_dt+B\sin\omega_dt\right),\qquad \omega_d=\omega_n\sqrt{1-\zeta^2}.
\end{equation}
\textbf{符号：} $x$ 为位移（m）；$t$ 为时间（s）；$\zeta$ 为阻尼比（无量纲）；$\omega_n,\omega_d$ 为无阻尼、阻尼后的角频率（rad/s）；$A,B$ 为由初始条件决定的长度常数（m）；$e$ 为自然对数底；$\cos,\sin$ 为三角函数。
\textbf{白话：} 正弦振荡外面套一个逐渐变小的包络。\textbf{条件与误用：} 外力为零且 $0\leq\zeta<1$；临界阻尼和过阻尼不能使用这条振荡表达式。数值 ODE 解可与本式比较验证程序。

\subsection{M14：线弹性与轴向杆变形}
\begin{equation}
\sigma=\frac{F}{A},\qquad \varepsilon=\frac{\Delta L}{L},\qquad \sigma=E\varepsilon,\qquad \Delta L=\frac{FL}{EA}.
\end{equation}
\textbf{符号：} $\sigma$ 为轴向正应力（Pa）；$F$ 为轴向拉力（N）；$A$ 为横截面积（m$^2$）；$\varepsilon$ 为轴向应变（无量纲）；$L$ 为原长（m）；$\Delta L$ 为伸长量（m）；$E$ 为杨氏模量（Pa）。
\textbf{白话：} 应力是每单位面积承担的力，应变是相对伸长比例。\textbf{条件与误用：} 均匀直杆、等截面、轴向载荷、小变形、线弹性范围；超出屈服后不能继续用直线关系；MPa 代入 SI 公式时要乘 $10^6$。

\subsection{M15：Euler--Bernoulli 梁与悬臂端挠度}
核对来源：\href{https://pastewka.github.io/Micromechanics/lecture/chapter04.html}{Lars Pastewka：Micromechanics，Beam kinematics 课程}。
\begin{equation}
EI_b\frac{d^4w}{dx^4}=q(x),\qquad I_b=\frac{bh^3}{12},\qquad \delta_{\rm tip}=\frac{FL^3}{3EI_b}.
\end{equation}
\textbf{符号：} $E$ 为杨氏模量（Pa）；$I_b$ 为截面对中性轴的面积二次矩（m$^4$，不是质量转动惯量）；$w$ 为与载荷正方向一致的横向挠度（m）；$x$ 为沿梁轴坐标（m）；$q(x)$ 为单位长度横向载荷（N/m）；$b$ 为矩形截面宽度（m）；$h$ 为弯曲方向高度（m）；$\delta_{\rm tip}$ 为悬臂自由端挠度大小（m）；$F$ 为自由端横向集中力（N）；$L$ 为梁长（m）。
\textbf{白话：} 梁越长越容易弯，高度增加尤其能增强抗弯能力。\textbf{条件与误用：} 细长、等刚度、小挠度梁；端挠度式仅用于一端固支、另一端受集中力；短厚梁或大转角要考虑剪切变形或几何非线性。通用四阶方程还需四个边界条件。

\section{光学：先分清光线模型与波动模型}
\subsection{O1：折射率与波长}
\begin{equation}
n=\frac{c_0}{v_p},\qquad \lambda=\frac{v_p}{f}=\frac{\lambda_0}{n}.
\end{equation}
\textbf{符号：} $n$ 为相折射率（无量纲）；$c_0=299792458\ \mathrm{m/s}$ 为真空光速；$v_p$ 为介质中相速度（m/s）；$\lambda,\lambda_0$ 为介质、真空波长（m）；$f$ 为光频率（Hz）。
\textbf{白话：} 光进入介质后频率通常不变，速度和波长改变。\textbf{条件与误用：} 透明介质中的单色波；色散时 $n$ 随频率变化。不要把群速度和相速度混用，也不要把 nm 直接当 m。

\subsection{O2：反射、折射与全反射}
\begin{equation}
\theta_r=\theta_i,\qquad n_1\sin\theta_i=n_2\sin\theta_t,\qquad \theta_c=\arcsin(n_2/n_1)\quad(n_1>n_2).
\end{equation}
\textbf{符号：} $\theta_i,\theta_r,\theta_t$ 为相对于界面法线的入射、反射、折射角（rad）；$n_1,n_2$ 为入射、透射介质折射率（无量纲）；$\theta_c$ 为全反射临界角（rad）；$\arcsin$ 为反正弦函数。
\textbf{白话：} 入射角从垂直界面的法线算；由高折射率介质射向低折射率介质且超过临界角时发生全反射。\textbf{条件与误用：} 平滑界面、几何光学近似；从界面本身量角会把正弦误写成余弦。来源：\href{https://openstax.org/books/university-physics-volume-3/pages/1-6-huygenss-principle}{OpenStax：惠更斯原理与反射折射}。

\subsection{O3：薄透镜成像}
\begin{equation}
\frac1f=\frac1u+\frac1v,\qquad M=\frac{h_i}{h_o}=-\frac vu.
\end{equation}
\textbf{符号：} $f$ 为焦距（m，会聚透镜正、发散透镜负）；$u$ 为物距（m，实物正）；$v$ 为像距（m，实像正、虚像负）；$M$ 为横向放大率（无量纲）；$h_i,h_o$ 为带方向的像高、物高（m）。
\textbf{白话：} 物体位置决定像出现在哪里以及是否倒立。\textbf{条件与误用：} 薄透镜、近轴小角度、采用此处符号约定；厚透镜需主平面，边缘光线需像差模型。\textbf{算例：} $f=0.10\ \mathrm m$，$u=0.30\ \mathrm m$，则 $v=0.15\ \mathrm m$、$M=-0.5$，得到倒立缩小实像。

\subsection{O4：光程与相位差}
\begin{equation}
\mathcal L=\int_C n(s)\,ds,\qquad \Delta\phi=\frac{2\pi}{\lambda_0}(\mathcal L_2-\mathcal L_1)+\phi_b.
\end{equation}
\textbf{符号：} $\mathcal L,\mathcal L_1,\mathcal L_2$ 为光程及两束光的光程（m）；$C$ 为光路；$s$ 为沿光路的弧长坐标（m）；$n(s)$ 为折射率（无量纲）；$ds$ 为微小路程（m）；$\Delta\phi$ 为两束光相位差（rad）；$\lambda_0$ 为真空波长（m）；$\phi_b$ 为反射等引入的附加相位差（rad）；$\pi$ 为圆周率。
\textbf{白话：} 光程把在不同介质里走路的“难易程度”换成等效长度。\textbf{条件与误用：} 相干单色波、约定同一比较点；薄膜反射可能多一个半波相位，不能只算几何距离。

\subsection{O5：两束光干涉强度}
\begin{equation}
I=I_1+I_2+2\sqrt{I_1I_2}\cos\Delta\phi.
\end{equation}
\textbf{符号：} $I$ 为合成光强（W/m$^2$）；$I_1,I_2$ 为两束光单独存在时的光强（W/m$^2$）；$\Delta\phi$ 为相位差（rad）。
\textbf{白话：} 两束光既能增强也能抵消，不能总把强度直接相加。\textbf{条件与误用：} 同频、相位关系稳定、偏振方向一致且空间重叠；非相干光长期平均后干涉项消失。测量噪声与部分相干会降低条纹对比度。

\subsection{O6：双缝条纹位置与间距}
\begin{equation}
d\sin\theta_m=m\lambda,\qquad y_m\simeq\frac{m\lambda D}{d},\qquad \Delta y\simeq\frac{\lambda D}{d}.
\end{equation}
\textbf{符号：} $d$ 为两缝中心间距（m）；$\theta_m$ 为第 $m$ 级亮纹相对中央方向的角度（rad）；$m$ 为整数级次，可为零或正负整数；$\lambda$ 为传播介质波长（m）；$y_m$ 为亮纹相对中央的屏上位置（m）；$D$ 为缝到屏距离（m）；$\Delta y$ 为相邻亮纹间距（m）；$\simeq$ 为近似相等。
\textbf{白话：} 波长越大、屏越远，条纹越疏。\textbf{条件与误用：} 相干窄缝，后两式另需远屏和小角度；$m=0$ 是中央亮纹；缝宽较大时还叠加单缝衍射包络。来源：\href{https://openstax.org/books/university-physics-volume-3/pages/3-1-youngs-double-slit-interference}{OpenStax：杨氏双缝干涉}。

\subsection{O7：单缝衍射}
\begin{equation}
I(\theta)=I_0\left(\frac{\sin\beta}{\beta}\right)^2,\qquad \beta=\frac{\pi a\sin\theta}{\lambda},\qquad a\sin\theta_m=m\lambda\quad(m=\pm1,\pm2,\ldots).
\end{equation}
\textbf{符号：} $I(\theta)$ 为方向 $\theta$ 上光强（W/m$^2$）；$I_0$ 为中央最大光强（W/m$^2$）；$\theta,\theta_m$ 为观察角、暗纹角（rad）；$\beta$ 为无量纲相位参数；$a$ 为缝宽（m）；$\lambda$ 为介质波长（m）；$m$ 为非零整数暗纹级次。
\textbf{白话：} 缝越窄，中央亮斑越宽。\textbf{条件与误用：} 均匀照明、远场单缝；中央 $\beta=0$ 处应取极限 1，程序直接除零会出错。暗纹公式不能把 $m=0$ 算成暗纹。来源：\href{https://openstax.org/books/university-physics-volume-3/pages/4-3-double-slit-diffraction}{OpenStax：衍射包络与双缝衍射}。

\subsection{O8：圆孔分辨极限}
\begin{equation}
\theta_{\min}\simeq1.22\frac{\lambda}{D_a}.
\end{equation}
\textbf{符号：} $\theta_{\min}$ 为 Rayleigh 判据的最小可分辨角（rad）；$\lambda$ 为介质波长（m）；$D_a$ 为圆形孔径直径（m）；1.22 为圆孔衍射第一零点相关系数。
\textbf{白话：} 更大的镜头口径原则上能分清更近的两个点。\textbf{条件与误用：} 理想圆孔、衍射受限成像、点源；像差、大气扰动、像素采样可能让实际分辨率更差；这不是所有图像算法的绝对分辨率上限。

\subsection{O9：Malus 偏振定律}
\begin{equation}
I=I_p\cos^2\theta,\qquad I_p=\frac12I_u\quad\text{（自然光通过理想起偏器）}.
\end{equation}
\textbf{符号：} $I$ 为检偏器后强度（W/m$^2$）；$I_p$ 为入射检偏器的线偏振强度（W/m$^2$）；$\theta$ 为入射偏振方向与检偏器透光轴夹角（rad）；$I_u$ 为起偏器前非偏振自然光强（W/m$^2$）。
\textbf{白话：} 检偏器只放过沿自身方向的光场分量。\textbf{条件与误用：} 理想偏振片；$\cos^2$ 用在线偏振光上，不是把自然光直接带入并任选一个角度。真实偏振片还存在透过率损失。

\subsection{O10：Beer--Lambert 衰减与吸光度}
核对来源：\href{https://goldbook.iupac.org/terms/view/B00626}{IUPAC Gold Book：Beer--Lambert 定律及单位约定}。
\begin{equation}
I=I_0e^{-\alpha\ell},\qquad A_{10}=\log_{10}\frac{I_0}{I}=\varepsilon c\ell.
\end{equation}
\textbf{符号：} $I,I_0$ 为透射、入射光强（W/m$^2$）；$\alpha$ 为自然指数衰减系数（m$^{-1}$）；$\ell$ 为光程长度（m）；$A_{10}$ 为十进制吸光度（无量纲）；$\varepsilon$ 为与十进制对数匹配的摩尔吸光系数（m$^2$/mol，若浓度用 mol/m$^3$）；$c$ 为浓度（mol/m$^3$）；$e$ 为自然对数底。
\textbf{白话：} 每走相同距离损失相同比例，形成指数衰减。\textbf{条件与误用：} 单色光、均匀介质、线性吸收；强散射、高浓度、化学平衡变化会偏离。仪器常用 L/(mol$\cdot$cm) 和 mol/L、cm，不能与 SI 数值直接混配；$\alpha$ 与 $\varepsilon c$ 相差 $\ln10$ 因子。

\section{热学：始终从能量收支开始}
\subsection{H1：热容量与显热}
\begin{equation}
Q=mc\,(T_2-T_1),\qquad C=mc.
\end{equation}
\textbf{符号：} $Q$ 为物体吸收的显热（J）；$m$ 为质量（kg）；$c$ 为相应过程的比热容（J/(kg$\cdot$K)）；$T_1,T_2$ 为初末温度（K 或同一摄氏温标）；$C$ 为总热容量（J/K）。
\textbf{白话：} 质量乘以每千克升温一度所需能量，再乘温差。\textbf{条件与误用：} 无相变、$c$ 近似常数；气体要区分定压比热和定容比热；水蒸发时不能只算升温。\textbf{算例：} 假设水的 $c=4200\ \mathrm{J/(kg\,K)}$，$m=1\ \mathrm{kg}$，升温 $10\ \mathrm K$，需 $42\ \mathrm{kJ}$。来源：\href{https://openstax.org/books/university-physics-volume-2/pages/1-4-heat-transfer-specific-heat-and-calorimetry}{OpenStax：比热与量热}。

\subsection{H2：热力学第一定律}
\begin{equation}
\Delta U=Q-W,\qquad W=\int_{V_1}^{V_2}p\,dV\quad\text{（准静态体积功）}.
\end{equation}
\textbf{符号：} $\Delta U$ 为系统内能增加量（J）；$Q$ 为外界传入系统的热量（J）；$W$ 为系统对外做功（J）；$p$ 为准静态过程中的边界压力（Pa）；$V,V_1,V_2$ 为体积及初末值（m$^3$）；$dV$ 为微小体积变化（m$^3$）。
\textbf{白话：} 流入的热，扣掉对外做的功，剩下的存成内能。\textbf{条件与误用：} 此写法用于封闭系统，忽略宏观动能、势能变化；开放管道还要计入流入流出焓与机械能；有些教材把对系统做功取正，使用时先统一符号。

\subsection{H3：理想气体状态方程}
\begin{equation}
pV=n_mR_uT,\qquad p=\rho R_sT,\qquad R_s=R_u/M_m.
\end{equation}
\textbf{符号：} $p$ 为绝对压力（Pa）；$V$ 为体积（m$^3$）；$n_m$ 为物质的量（mol）；$R_u\approx8.31446\ \mathrm{J/(mol\,K)}$ 为通用气体常数；$T$ 为绝对温度（K）；$\rho$ 为密度（kg/m$^3$）；$R_s$ 为比气体常数（J/(kg$\cdot$K)）；$M_m$ 为摩尔质量（kg/mol）。
\textbf{白话：} 同种气体压得更密或加热，压力通常升高。\textbf{条件与误用：} 稀薄、远离凝结的气体近似；不能用表压或摄氏温度；高压或接近临界点时要用真实气体状态方程。

\subsection{H4：Fourier 导热定律}
\begin{equation}
\mathbf q''=-k\nabla T,\qquad \dot Q=\frac{kA}{L}(T_h-T_c)\quad\text{（稳态平板）}.
\end{equation}
\textbf{符号：} $\mathbf q''$ 为热流密度向量（W/m$^2$）；$k$ 为导热系数（W/(m$\cdot$K)）；$\nabla T$ 为温度空间梯度（K/m）；$\dot Q$ 为从热侧到冷侧的传热率（W）；$A$ 为垂直热流的面积（m$^2$）；$L$ 为厚度（m）；$T_h,T_c$ 为两面温度（K 或同一摄氏温标）。
\textbf{白话：} 热往温度低的地方走，负号保证方向正确。\textbf{条件与误用：} 第一式为各向同性介质；第二式另需一维、稳态、无内热源、常导热系数。不要把热流密度与总功率混淆。来源：\href{https://openstax.org/books/university-physics-volume-2/pages/1-6-mechanisms-of-heat-transfer}{OpenStax：传热机制}。

\subsection{H5：Newton 对流换热}
\begin{equation}
\dot Q=hA(T_s-T_\infty).
\end{equation}
\textbf{符号：} $\dot Q$ 为表面传向流体的热功率（W）；$h$ 为对流换热系数（W/(m$^2\cdot$K)）；$A$ 为换热面积（m$^2$）；$T_s$ 为表面温度（K）；$T_\infty$ 为远离表面的流体主体温度（K）。
\textbf{白话：} 表面比空气热得越多、面积越大，通常散热越快。\textbf{条件与误用：} $h$ 是由流动、几何、物性决定的经验或计算量；不是材料固有常数。物体内部温度不均匀时，不能把中心温度直接当表面温度。

\subsection{H6：Stefan--Boltzmann 热辐射}
\begin{equation}
\dot Q_{\rm rad}=\epsilon\sigma_{\rm SB}A(T_s^4-T_{\rm sur}^4).
\end{equation}
\textbf{符号：} $\dot Q_{\rm rad}$ 为净向外辐射功率（W）；$\epsilon$ 为表面有效发射率（无量纲，0 到 1）；$\sigma_{\rm SB}\approx5.670374419\times10^{-8}\ \mathrm{W/(m^2\,K^4)}$ 为 Stefan--Boltzmann 常数；$A$ 为表面积（m$^2$）；$T_s,T_{\rm sur}$ 为表面和大包围环境绝对温度（K）。
\textbf{白话：} 高温物体辐射增长非常快，因为温度被四次方。\textbf{条件与误用：} 小灰体位于温度均匀的大包围环境，且可用该有效发射率近似；多表面腔体要考虑视角因子和反射；绝不能把摄氏温度四次方。

\subsection{H7：平板热阻、对流热阻与串联}
\begin{equation}
R_{\rm cond}=\frac{L}{kA},\quad R_{\rm conv}=\frac1{hA},\quad R_{\rm total}=\sum_jR_j,\quad \dot Q=\frac{T_h-T_c}{R_{\rm total}}.
\end{equation}
\textbf{符号：} $R_{\rm cond},R_{\rm conv},R_j,R_{\rm total}$ 为导热、对流、第 $j$ 段及总热阻（K/W）；$L$ 为层厚（m）；$k$ 为导热系数（W/(m$\cdot$K)）；$A$ 为对应面积（m$^2$）；$h$ 为换热系数（W/(m$^2\cdot$K)）；$j$ 为层或边界编号；$T_h,T_c$ 为热冷端温度（K）；$\dot Q$ 为热功率（W）。
\textbf{白话：} 把传热路径像电路一样一段段相加。\textbf{条件与误用：} 稳态、同一热流串联；并行传热路径应加热导而不是热阻；接触不良要另加接触热阻。

\subsection{H8：圆筒壁导热热阻}
\begin{equation}
R_{\rm cyl}=\frac{\ln(r_o/r_i)}{2\pi kL},\qquad \dot Q=\frac{T_i-T_o}{R_{\rm cyl}}.
\end{equation}
\textbf{符号：} $R_{\rm cyl}$ 为圆筒壁热阻（K/W）；$r_i,r_o$ 为内外半径（m）；$k$ 为导热系数（W/(m$\cdot$K)）；$L$ 为轴向长度（m）；$T_i,T_o$ 为内外壁温度（K）；$\dot Q$ 为向外传热功率（W）；$\ln$ 为自然对数。
\textbf{白话：} 圆筒越向外面积越大，所以不能随便拿内表面积套平板式。\textbf{条件与误用：} 轴向足够长、一维径向稳态、无热源、常物性；$L$ 是管长，不是壁厚。

\subsection{H9：集总热容冷却与时间常数}
\begin{equation}
mc\frac{dT}{dt}=-hA(T-T_\infty),\quad T(t)=T_\infty+(T_0-T_\infty)e^{-t/\tau},\quad \tau=\frac{mc}{hA}.
\end{equation}
\textbf{符号：} $m$ 为物体质量（kg）；$c$ 为比热容（J/(kg$\cdot$K)）；$T,T_0$ 为均匀物体温度及初值（K）；$t$ 为时间（s）；$h$ 为换热系数（W/(m$^2\cdot$K)）；$A$ 为外表面积（m$^2$）；$T_\infty$ 为恒定环境温度（K）；$\tau$ 为时间常数（s）；$e$ 为自然对数底。
\textbf{白话：} 经过一个时间常数，物体与环境的温差剩下原来的约 36.8\%。\textbf{条件与误用：} 内部近似等温、无内热源、常物性、仅上述对流机制；加入明显辐射后不能仍直接用本指数解。

\subsection{H10：Biot 数与 Fourier 数}
\begin{equation}
\mathrm{Bi}=\frac{hL_c}{k_s},\qquad L_c=\frac VA,\qquad \mathrm{Fo}=\frac{\alpha t}{L_d^2},\qquad \alpha=\frac{k_s}{\rho c}.
\end{equation}
\textbf{符号：} $\mathrm{Bi}$ 为内外传热阻力比较的 Biot 数（无量纲）；$h$ 为表面对流系数（W/(m$^2\cdot$K)）；$L_c$ 为集总判据特征长度（m）；$V$ 为物体体积（m$^3$）；$A$ 为参与换热的面积（m$^2$）；$k_s$ 为固体导热系数（W/(m$\cdot$K)）；$\mathrm{Fo}$ 为热扩散时间尺度比（无量纲）；$\alpha$ 为热扩散率（m$^2$/s）；$t$ 为时间（s）；$L_d$ 为所研究扩散距离（m）；$\rho$ 为密度（kg/m$^3$）；$c$ 为比热容（J/(kg$\cdot$K)）。
\textbf{白话：} Bi 小表示内部传热快，适合把整个物体当一个温度；Fo 衡量热扩散走了多远。\textbf{条件与误用：} 常用 $\mathrm{Bi}<0.1$ 是工程近似筛选而非普适严格误差保证；不同解析解的特征长度定义必须核对，不能把板厚与半厚混用。

\subsection{H11：含热源的导热方程}
\begin{equation}
\rho c\frac{\partial T}{\partial t}=\nabla\cdot(k\nabla T)+q_v.
\end{equation}
\textbf{符号：} $\rho$ 为密度（kg/m$^3$）；$c$ 为比热容（J/(kg$\cdot$K)）；$T$ 为温度场（K）；$t$ 为时间（s）；$k$ 为导热系数（W/(m$\cdot$K)）；$q_v$ 为单位体积产热率（W/m$^3$）；$\partial$ 为偏导数；$\nabla T$ 为梯度（K/m）；$\nabla\cdot$ 为散度，计算热流在局部的收支。
\textbf{白话：} 一小块材料升温，是因为流入的热比流出的多，或它自己在发热。\textbf{条件与误用：} 静止介质、适当定义热容；流动流体还要加入对流输运项；变 $k$ 时不能把 $k$ 直接移出散度。初始条件与边界条件缺一不可。

\subsection{H12：相变潜热与界面 Stefan 条件}
拓展来源：\href{https://cea-lbm-saclay.github.io/LBM_Saclay_Documentation/src_doc/05_COURSES/Thermo-PhaseTransition/Phase_Change_Models.html}{CEA Saclay：固液相变模型课程}；本节采用固体在左的坐标约定，比较文献时须先统一相区方向。
\begin{equation}
Q_{\rm lat}=mL_f,\qquad \rho L_f\frac{ds}{dt}=k_s\left.\frac{\partial T_s}{\partial x}\right|_{s^-}-k_l\left.\frac{\partial T_l}{\partial x}\right|_{s^+}.
\end{equation}
\textbf{符号：} $Q_{\rm lat}$ 为相变热量（J）；$m$ 为发生相变的质量（kg）；$L_f$ 为单位质量熔化潜热（J/kg）；$\rho$ 为共同近似密度（kg/m$^3$）；$x$ 为向右坐标（m）；$s(t)$ 为界面位置（m），固体在左、液体在右，正 $ds/dt$ 表示凝固推进；$t$ 为时间（s）；$k_s,k_l$ 为固、液导热系数（W/(m$\cdot$K)）；$T_s,T_l$ 为固、液温度场（K）；$s^-,s^+$ 表示从固、液侧趋近界面。
\textbf{白话：} 凝固释放的潜热等于界面两边热流的净差。\textbf{条件与误用：} 平面尖锐界面、无界面储热、等密度近似，并在界面给出熔点条件；方向定义变化会改变符号。不要将辐射常数 Stefan--Boltzmann 定律与相变 Stefan 条件混为一谈。

\section{流体：先看守恒，再看黏性与尺度}
\subsection{F1：密度、压力与静水压力}
\begin{equation}
\rho=\frac mV,\qquad p=\frac{F_n}{A},\qquad p(z)=p_0+\rho gz.
\end{equation}
\textbf{符号：} $\rho$ 为密度（kg/m$^3$）；$m$ 为质量（kg）；$V$ 为体积（m$^3$）；$p$ 为均匀法向压力（Pa）；$F_n$ 为法向压力合力（N）；$A$ 为受压面积（m$^2$）；$p_0$ 为液面压力（Pa）；$g$ 为重力加速度（m/s$^2$）；$z$ 为从液面向下量的深度（m）。
\textbf{白话：} 同一静止液体越深压力越大。\textbf{条件与误用：} 静止、密度近似不变；静水式这里向下为正，不要与向上坐标公式混用；表压与绝对压力需标明。

\subsection{F2：浮力}
\begin{equation}
F_b=\rho_f gV_{\rm disp}.
\end{equation}
\textbf{符号：} $F_b$ 为向上浮力大小（N）；$\rho_f$ 为周围流体密度（kg/m$^3$）；$g$ 为重力加速度（m/s$^2$）；$V_{\rm disp}$ 为排开流体体积（m$^3$），即浸没部分体积。
\textbf{白话：} 浮力等于被挤开的那部分流体的重量。\textbf{条件与误用：} 静水、均匀密度近似；漂浮时不能用物体全部体积；水底完全密封接触等特殊边界不能直接照套。

\subsection{F3：流量与连续性方程}
\begin{equation}
Q_v=\int_A\mathbf u\cdot\mathbf n\,dA,\quad \dot m=\rho Q_v,\quad \frac{\partial\rho}{\partial t}+\nabla\cdot(\rho\mathbf u)=0,\quad A_1\bar u_1=A_2\bar u_2.
\end{equation}
\textbf{符号：} $Q_v$ 为体积流量（m$^3$/s）；$A$ 为截面，$dA$ 为微小面积（m$^2$）；$\mathbf u$ 为速度场（m/s）；$\mathbf n$ 为截面单位法向量（无量纲）；$\dot m$ 为质量流量（kg/s）；$\rho$ 为密度（kg/m$^3$）；$t$ 为时间（s）；$\nabla\cdot$ 为散度；$A_1,A_2$ 为两截面积（m$^2$）；$\bar u_1,\bar u_2$ 为两截面的法向平均流速（m/s）。
\textbf{白话：} 不漏水、不储水的管道，进来多少就出去多少；变窄处平均流速会增大。\textbf{条件与误用：} $\dot m=\rho Q_v$ 需截面密度均匀；最后一式需稳态不可压且无支流；中心最大速度不能代替平均速度。

\subsection{F4：不可压 Newton 流体的 Navier--Stokes 方程}
\begin{equation}
\rho\left(\frac{\partial\mathbf u}{\partial t}+(\mathbf u\cdot\nabla)\mathbf u\right)=-\nabla p+\mu\nabla^2\mathbf u+\rho\mathbf g,\qquad \nabla\cdot\mathbf u=0.
\end{equation}
\textbf{符号：} $\rho$ 为恒定密度（kg/m$^3$）；$\mathbf u$ 为速度（m/s）；$t$ 为时间（s）；$p$ 为压力（Pa）；$\mu$ 为恒定动力黏度（Pa$\cdot$s）；$\mathbf g$ 为重力加速度向量（m/s$^2$）；$\nabla$ 为梯度算子；$\nabla^2$ 为对各空间方向二阶导数求和；$(\mathbf u\cdot\nabla)\mathbf u$ 为流体运动到不同位置导致的加速度（m/s$^2$）。
\textbf{白话：} 流体版牛顿第二定律：惯性等于压力作用、黏性作用与重力。\textbf{条件与误用：} 连续介质、不可压、Newton 流体、常黏度；可压流和非 Newton 流体要改方程。必须给初边值条件；直接离散并不自动保证压力速度耦合正确。来源：\href{https://www.grc.nasa.gov/www/k-12/airplane/nseqs.html}{NASA Glenn：Navier--Stokes 方程概述}。

\subsection{F5：Bernoulli 能量式}
\begin{equation}
\frac p{\rho g}+\frac{u^2}{2g}+z=H.
\end{equation}
\textbf{符号：} $p$ 为压力（Pa）；$\rho$ 为流体密度（kg/m$^3$）；$g$ 为重力加速度（m/s$^2$）；$u$ 为局部速率（m/s）；$z$ 为向上的高程（m）；$H$ 为沿同一流线的总水头常数（m）。
\textbf{白话：} 压力、高度和速度可以互相交换能量。\textbf{条件与误用：} 稳态、不可压、忽略黏性损耗、无泵或涡轮做功，通常只沿同一流线成立；存在明显管摩擦时不能硬设两点 $H$ 相等。三个加项的单位都是长度。

\subsection{F6：有泵和损失的管道能量式}
\begin{equation}
\frac{p_1}{\rho g}+\alpha_1\frac{\bar u_1^2}{2g}+z_1+h_p=\frac{p_2}{\rho g}+\alpha_2\frac{\bar u_2^2}{2g}+z_2+h_t+h_L.
\end{equation}
\textbf{符号：} $p_1,p_2$ 为两截面压力（Pa）；$\rho$ 为恒定密度（kg/m$^3$）；$g$ 为重力加速度（m/s$^2$）；$\bar u_1,\bar u_2$ 为截面平均速度（m/s）；$\alpha_1,\alpha_2$ 为动能修正系数（无量纲）；$z_1,z_2$ 为高程（m）；$h_p$ 为泵加入水头（m）；$h_t$ 为涡轮取出水头（m）；$h_L$ 为不可逆损失水头（m）。
\textbf{白话：} 泵加能量，涡轮和阻力拿走能量。\textbf{条件与误用：} 稳态不可压；均匀剖面修正系数为 1，充分发展圆管层流为 2；不能把泵输入电功率全部当流体有效功率。

\subsection{F7：Darcy--Weisbach 沿程损失与局部损失}
\begin{equation}
\Delta p_f=f_D\frac LD\frac{\rho\bar u^2}{2},\qquad \Delta p_{\rm local}=K\frac{\rho\bar u^2}{2},\qquad h_L=\frac{\Delta p_{\rm loss}}{\rho g}.
\end{equation}
\textbf{符号：} $\Delta p_f,\Delta p_{\rm local},\Delta p_{\rm loss}$ 为沿程、局部、总损失对应压差（Pa）；$f_D$ 为 Darcy 摩擦因子（无量纲）；$L$ 为管长（m）；$D$ 为圆管内径（m）；$\rho$ 为密度（kg/m$^3$）；$\bar u$ 为平均速度（m/s）；$K$ 为按指定参考速度定义的局部阻力系数（无量纲）；$h_L$ 为总损失水头（m）；$g$ 为重力加速度（m/s$^2$）。
\textbf{白话：} 长管损失按长度累积，弯头阀门另外计。\textbf{条件与误用：} 管道经验模型；Darcy 因子是 Fanning 因子的 4 倍，必须查清文献定义。变径附件的 $K$ 必须配套其规定的截面速度。

\subsection{F8：圆管 Poiseuille 层流}
\begin{equation}
Q_v=\frac{\pi R^4\Delta p}{8\mu L},\qquad \bar u=\frac{Q_v}{\pi R^2},\qquad f_D=\frac{64}{\mathrm{Re}}.
\end{equation}
\textbf{符号：} $Q_v$ 为体积流量（m$^3$/s）；$R$ 为管内半径（m）；$\Delta p$ 为水平直管两端压差（Pa）；$\mu$ 为动力黏度（Pa$\cdot$s）；$L$ 为管长（m）；$\bar u$ 为平均速度（m/s）；$f_D$ 为 Darcy 摩擦因子（无量纲）；$\mathrm{Re}$ 为以管径和平均速度定义的 Reynolds 数（无量纲）。
\textbf{白话：} 同样压差下，半径翻倍会让该模型的流量变为 16 倍。\textbf{条件与误用：} 长直圆管、稳态、不可压 Newton 流体、充分发展层流、无滑移；入口区或湍流不能套；半径与直径弄错会引入 16 倍误差。来源：\href{https://openstax.org/books/university-physics-volume-1/pages/14-7-viscosity-and-turbulence}{OpenStax：黏性、Poiseuille 定律与湍流}。

\subsection{F9：Reynolds 数与运动黏度}
\begin{equation}
\mathrm{Re}=\frac{\rho UL}{\mu}=\frac{UL}{\nu},\qquad \nu=\frac\mu\rho.
\end{equation}
\textbf{符号：} $\mathrm{Re}$ 为惯性与黏性作用之比的 Reynolds 数（无量纲）；$\rho$ 为密度（kg/m$^3$）；$U$ 为特征速度（m/s）；$L$ 为特征长度（m）；$\mu$ 为动力黏度（Pa$\cdot$s）；$\nu$ 为运动黏度（m$^2$/s）。
\textbf{白话：} 数值越小，黏性越能把流动“抹平”。\textbf{条件与误用：} $U,L$ 按问题定义；圆管常取平均速度与直径，约 2300 以下通常为层流筛选、2300--4000 为过渡区，但扰动和入口会改变转捩。不能把这组管流阈值当平板、球体的通用阈值。来源：\href{https://www.grc.nasa.gov/WWW/k-12/airplane/reynolds.html}{NASA Glenn：Reynolds 数}。

\subsection{F10：Mach 数与声速}
\begin{equation}
\mathrm{Ma}=\frac Ua,\qquad a=\sqrt{\gamma R_sT}.
\end{equation}
\textbf{符号：} $\mathrm{Ma}$ 为 Mach 数（无量纲）；$U$ 为流速（m/s）；$a$ 为当地声速（m/s）；$\gamma$ 为定压与定容比热之比（无量纲）；$R_s$ 为该气体比气体常数（J/(kg$\cdot$K)）；$T$ 为绝对温度（K）。
\textbf{白话：} Mach 数比较流动快慢与压力扰动传播快慢。\textbf{条件与误用：} 声速式用于理想气体的局部等熵小扰动；$\mathrm{Ma}<0.3$ 常作为弱压缩性初筛，但强加热等仍可能造成明显密度变化。来源：\href{https://www1.grc.nasa.gov/beginners-guide-to-aeronautics/mass-flow-rate-equations/}{NASA Glenn：质量流率、Mach 数与声速}。

\subsection{F11：二次阻力与低 Reynolds 球体阻力}
\begin{equation}
F_D=\frac12C_D\rho A_pU^2,\qquad F_{\rm Stokes}=6\pi\mu RU.
\end{equation}
\textbf{符号：} $F_D,F_{\rm Stokes}$ 为阻力大小（N）；$C_D$ 为与形状和 Reynolds 数等有关的阻力系数（无量纲）；$\rho$ 为流体密度（kg/m$^3$）；$A_p$ 为与系数定义匹配的参考投影面积（m$^2$）；$U$ 为物体相对流体速率（m/s）；$\mu$ 为动力黏度（Pa$\cdot$s）；$R$ 为球半径（m）。
\textbf{白话：} 阻力方向反对相对运动；非常慢的小球由黏性主导，阻力与速度一次方相关。\textbf{条件与误用：} Stokes 式需孤立刚性球、无滑移、近似无限流体、$\mathrm{Re}\ll1$；二次式的 $C_D$ 一般并非恒定。注意减去流体背景速度后才是相对速度。

\subsection{F12：Nusselt、Prandtl 与 P\'eclet 数}
核对来源：\href{https://bpb-us-e1.wpmucdn.com/sites.mit.edu/dist/6/1810/files/2024/04/AHTTv600.pdf}{Lienhard 与 Lienhard：A Heat Transfer Textbook，第六版，符号表及对流换热章节}。
\begin{equation}
\mathrm{Nu}=\frac{hL}{k_f},\qquad \mathrm{Pr}=\frac{\mu c_p}{k_f}=\frac\nu\alpha,\qquad \mathrm{Pe}=\mathrm{Re}\,\mathrm{Pr}=\frac{UL}{\alpha}.
\end{equation}
\textbf{符号：} $\mathrm{Nu}$ 为对流与跨特征长度导热强弱之比（无量纲）；$h$ 为换热系数（W/(m$^2\cdot$K)）；$L$ 为特征长度（m）；$k_f$ 为流体导热系数（W/(m$\cdot$K)）；$\mathrm{Pr}$ 为动量扩散与热扩散之比（无量纲）；$\mu$ 为动力黏度（Pa$\cdot$s）；$c_p$ 为定压比热（J/(kg$\cdot$K)）；$\nu$ 为运动黏度（m$^2$/s）；$\alpha$ 为流体热扩散率（m$^2$/s）；$\mathrm{Pe}$ 为输运与热扩散之比（无量纲）；$\mathrm{Re}$ 为采用同一速度长度定义的 Reynolds 数；$U$ 为特征速度（m/s）。
\textbf{白话：} 无量纲数把不同尺寸的实验放在同一尺度下比较；知道 Nu 就能反求 $h$。\textbf{条件与误用：} 查换热关联式时必须匹配几何、流态、边界条件及物性评估温度；不要只因名字相同就互换特征长度。

\subsection{F13：流动携热与泵有效功率}
\begin{equation}
\dot Q=\dot m c_p(T_{\rm out}-T_{\rm in}),\qquad P_{\rm hyd}=\Delta p\,Q_v,\qquad P_{\rm in}=\frac{P_{\rm hyd}}\eta.
\end{equation}
\textbf{符号：} $\dot Q$ 为流体吸收热功率（W）；$\dot m$ 为质量流量（kg/s）；$c_p$ 为定压比热（J/(kg$\cdot$K)）；$T_{\rm out},T_{\rm in}$ 为出口、入口混合平均温度（K）；$P_{\rm hyd}$ 为液压有效功率（W）；$\Delta p$ 为泵提供的等效总压升（Pa）；$Q_v$ 为体积流量（m$^3$/s）；$P_{\rm in}$ 为输入功率（W）；$\eta$ 为输入到液压功率的总效率（无量纲，0 到 1）。
\textbf{白话：} 流量越大，同样升温需要的热越多；泵输入还要补上效率损失。\textbf{条件与误用：} 热量式需稳态、无相变、常比热、机械能变化可忽略；功率式用于不可压泵流，若只量静压需补偿速度与高程差。

\section{怎样把物理工具箱接到比赛模型}
\begin{enumerate}
\item 写研究对象：例如“长度为某值的加热板”“阀门之间的一段管道”，避免整个工程系统一次全建。
\item 写假设并说明证据：是否一维、是否稳态、是否小变形、是否可忽略压缩性。数量级分析比“为了简化”更有说服力。
\item 列参数表：名称、符号、单位、数值、来源、是否待估。材料参数随温度变化时应给适用范围。
\item 先写守恒式，再代入本构式。守恒式管收支，本构式说明力与变形、热流与温差、黏性与速度梯度如何联系。
\item 设置初始和边界条件，选数值算法，检查量纲、守恒误差、极限情况和网格收敛。
\item 用实验或题目数据校准，再用未参与拟合的数据验证。不能以曲线看起来平滑作为模型正确的证据。
\end{enumerate}

\chapter{概率、统计与不确定性：结果到底有多可信}
本章把“一个数算出来了”升级成“这个数在什么假设下可信”。大写 $X$ 通常表示尚未观测的随机变量，小写 $x$ 表示已经看到的具体数值；帽子 $\hat{\theta}$ 表示由数据估计的参数；横线表示平均。概率无单位，概率密度的单位则是随机变量单位的倒数。以下公式中的求和符号，意思是把指定范围的项目逐个加起来。

\section{概率的基本语言}
\subsection{S01：条件概率与乘法公式}
\begin{equation}
 P(A\mid B)=\frac{P(A\cap B)}{P(B)},\qquad P(A\cap B)=P(A\mid B)P(B).
\end{equation}
\textbf{人话：}已经知道 $B$ 发生后，只在 $B$ 的范围里重新计算 $A$ 发生的比例。例如只看高温工况下设备失效的比例。
\textbf{符号：}$A,B$ 是事件；$P$ 表示概率；竖线表示“已知”；$\cap$ 表示同时发生。所有概率无单位。
\textbf{条件与误用：}要求 $P(B)>0$。通常 $P(A\mid B)$ 不等于 $P(B\mid A)$；“失效设备常常高温”不等于“高温设备常常失效”。独立时条件概率才等于原来的 $P(A)$。基础阅读：\url{https://ocw.mit.edu/courses/6-041sc-probabilistic-systems-analysis-and-applied-probability-fall-2013/}。

\subsection{S02：全概率与贝叶斯公式}
\begin{equation}
 P(B)=\sum_{j=1}^{m}P(B\mid A_j)P(A_j),\qquad
 P(A_i\mid B)=\frac{P(B\mid A_i)P(A_i)}{\sum_{j=1}^{m}P(B\mid A_j)P(A_j)}.
\end{equation}
\textbf{人话：}先把总体分成互不重叠的来源，再按各来源占比加权。贝叶斯公式用新证据更新来源的可能性。
\textbf{符号：}$A_1,\ldots,A_m$ 是覆盖全部可能情况且互斥的 $m$ 类事件；$i$ 是想追查的类别；$j$ 是求和用的临时下标；$B$ 是已看到的证据；$P(A_i)$ 是先验，$P(A_i\mid B)$ 是后验，$P(B\mid A_i)$ 是似然意义上的证据概率。概率均无单位。
\textbf{条件与误用：}分母必须大于零；各类必须没有重复、没有遗漏。不能把先验随便设为相等却不说明依据；故障率很低时，即使报警器准确，报警后也未必很可能真故障。基础阅读：\url{https://ocw.mit.edu/courses/6-041sc-probabilistic-systems-analysis-and-applied-probability-fall-2013/}。

\subsection{S03：期望——按可能性加权的平均}
\begin{equation}
 \mathrm{E}[X]=\sum_k x_kp_k\quad\text{或}\quad
 \mathrm{E}[X]=\int_{-\infty}^{\infty}x f_X(x)\,\mathrm{d}x.
\end{equation}
\textbf{人话：}结果越常出现，在平均中占的比重越大。它描述长期平均，不保证单次结果接近这个数。
\textbf{符号：}$x_k$ 是离散取值，$p_k=P(X=x_k)$；$k$ 枚举全部取值；$f_X$ 是连续变量的概率密度；积分是连续情形的加权求和；$\mathrm{d}x$ 是积分变量的微小区间。$X$ 若以米计，期望也以米计，密度以 $\mathrm{m}^{-1}$ 计。
\textbf{条件与误用：}通常要求绝对可积，即绝对值的平均有限。不要把连续密度 $f_X(x)$ 当作点概率，连续变量单点概率通常为零。基础阅读：\url{https://ocw.mit.edu/courses/6-041sc-probabilistic-systems-analysis-and-applied-probability-fall-2013/}。

\subsection{S04：方差与标准差——波动有多大}
\begin{equation}
 \operatorname{Var}(X)=\mathrm{E}[(X-\mu)^2]=\mathrm{E}[X^2]-\mu^2,
 \qquad \sigma=\sqrt{\operatorname{Var}(X)},\quad \mu=\mathrm{E}[X].
\end{equation}
\textbf{人话：}先计算离平均值有多远，再平方平均，避免正负抵消；开平方后回到原单位。
\textbf{符号：}$\mu$ 是总体均值；$\sigma$ 是总体标准差；$\mathrm{E}$ 是期望。若 $X$ 单位为 K，方差为 $\mathrm{K}^2$，标准差为 K。
\textbf{条件与误用：}需要二阶矩有限。均值相同的设备稳定性仍可能完全不同；标准差描述单个观测的波动，均值的标准误描述估计平均值的不确定性，二者不同。基础阅读：\url{https://ocw.mit.edu/courses/6-041sc-probabilistic-systems-analysis-and-applied-probability-fall-2013/}。

\subsection{S05：协方差与相关系数}
\begin{equation}
 \operatorname{Cov}(X,Y)=\mathrm{E}[(X-\mu_X)(Y-\mu_Y)],\qquad
 \rho_{XY}=\frac{\operatorname{Cov}(X,Y)}{\sigma_X\sigma_Y}.
\end{equation}
\textbf{人话：}两个量是否经常一起偏高或一起偏低？相关系数把这种线性同步变化压到 $[-1,1]$ 范围。
\textbf{符号：}$X,Y$ 是两个变量；$\mu_X,\mu_Y$ 是各自均值；$\sigma_X,\sigma_Y$ 是各自标准差；下标 $XY$ 表示这两个变量之间的关系。协方差单位为两个变量单位的乘积，$\rho_{XY}$ 无单位。
\textbf{条件与误用：}二阶矩有限且两个标准差非零。零相关只排除线性关系，不能一般推出独立；相关不证明因果。温度和能耗共同上升，可能由负荷这个第三变量驱动。基础阅读：\url{https://ocw.mit.edu/courses/6-041sc-probabilistic-systems-analysis-and-applied-probability-fall-2013/}。

\section{先根据产生机制选择分布}
\subsection{S06：伯努利与二项分布}
\begin{equation}
 P(X=k)=\binom{n}{k}p^k(1-p)^{n-k},\quad k=0,1,\ldots,n;
 \quad \mathrm{E}[X]=np,\quad\operatorname{Var}(X)=np(1-p).
\end{equation}
\textbf{人话：}独立重复 $n$ 次成功或失败的试验，统计成功次数。$n=1$ 就是伯努利分布。
\textbf{符号：}$X,k$ 为成功次数；$n$ 为总次数；$p$ 为每次成功概率；$\binom{n}{k}=n!/[k!(n-k)!]$ 为从 $n$ 次中选 $k$ 次成功位置的方法数，感叹号表示阶乘，$0!=1$。次数与概率按无量纲处理。
\textbf{条件与误用：}次数固定，各次独立且成功概率相同。若设备之间受共同电源故障影响，就不应直接假设独立。定义核对：\url{https://docs.scipy.org/doc/scipy/reference/generated/scipy.stats.binom.html}。

\subsection{S07：泊松分布——固定区间内的事件次数}
\begin{equation}
 P(N=k)=\frac{\lambda^k e^{-\lambda}}{k!},\quad k=0,1,2,\ldots;
 \qquad \mathrm{E}[N]=\operatorname{Var}(N)=\lambda=rt.
\end{equation}
\textbf{人话：}描述固定时间内随机到达的数量，例如稳定工况下计数器记录的稀疏事件。
\textbf{符号：}$N$ 是次数；$k$ 是具体次数；$r$ 为平均发生率，单位如 $\mathrm{s}^{-1}$；$t$ 是区间长度，单位 s；$\lambda$ 是区间内平均次数，无量纲；$e$ 为自然常数，约 $2.71828$。
\textbf{条件与误用：}齐次泊松过程要求不重叠区间增量独立、发生率恒定、极短时间多次事件概率足够小。明显成群发生或方差远大于均值时，应检查过度离散。定义核对：\url{https://docs.scipy.org/doc/scipy/reference/generated/scipy.stats.poisson.html}。

\subsection{S08：正态分布——对称的随机波动}
\begin{equation}
 f_X(x)=\frac{1}{\sigma\sqrt{2\pi}}\exp\!\left[-\frac{(x-\mu)^2}{2\sigma^2}\right],\qquad Z=\frac{X-\mu}{\sigma}.
\end{equation}
\textbf{人话：}钟形曲线，中心由均值决定、胖瘦由标准差决定；标准化把不同尺度转为“离均值几个标准差”。
\textbf{符号：}$x$ 是取值，$\mu$ 是均值，$\sigma>0$ 是标准差；$\pi$ 是圆周率；$\exp(a)=e^a$；$Z$ 是无量纲标准正态变量。$x,\mu,\sigma$ 的单位相同，密度单位相反。
\textbf{条件与误用：}正态是分布假设，需要看残差或机制。绝对不能只因画出了平滑曲线就认定正态；正态分布允许负值，不适合直接描述跨度很大的严格正量。定义核对：\url{https://docs.scipy.org/doc/scipy/reference/generated/scipy.stats.norm.html}。

\subsection{S09：均匀分布——只有上下界信息时的明确假设}
\begin{equation}
 f_X(x)=\frac{1}{b-a}\ (a\le x\le b),\quad f_X(x)=0\ \text{（区间外）};
 \quad\mathrm{E}[X]=\frac{a+b}{2},\quad\operatorname{Var}(X)=\frac{(b-a)^2}{12}.
\end{equation}
\textbf{人话：}相同长度的小区间被抽到的机会相同。
\textbf{符号：}$a$ 为下界，$b$ 为上界，要求 $a<b$；$x$ 是可能值。$a,b,x$ 同单位，方差是其单位平方。
\textbf{条件与误用：}只知道范围并不逻辑上推出均匀分布。用它模拟参数不确定性时，要写明这是一项假设，并检查改用其他分布后结论是否改变。定义核对：\url{https://docs.scipy.org/doc/scipy/reference/generated/scipy.stats.uniform.html}。

\subsection{S10：指数分布——等待时间}
\begin{equation}
 f_T(t)=r e^{-rt}\ (t\ge0),\qquad P(T>t)=e^{-rt},\qquad
 \mathrm{E}[T]=1/r,\quad\operatorname{Var}(T)=1/r^2.
\end{equation}
\textbf{人话：}稳定泊松到达过程中，下一次事件还要等多久。
\textbf{符号：}$T$ 是非负等待时间，$t$ 是指定时间，单位 s；$r>0$ 是发生率，单位 $\mathrm{s}^{-1}$；$P(T>t)$ 是等了 $t$ 秒还没发生的概率。
\textbf{条件与误用：}它具有无记忆性，相当于事件风险不随已等待时间变化。磨损越来越严重的寿命数据通常不满足这个条件。SciPy 的尺度参数为 $1/r$，不要把率直接当尺度。定义核对：\url{https://docs.scipy.org/doc/scipy/reference/generated/scipy.stats.expon.html}。

\section{从有限数据估计总体}
\subsection{S11：样本均值、样本方差与标准误}
\begin{equation}
 \bar{x}=\frac1n\sum_{i=1}^{n}x_i,\quad
 s^2=\frac1{n-1}\sum_{i=1}^{n}(x_i-\bar{x})^2,\quad
 s=\sqrt{s^2},\quad\widehat{\operatorname{SE}}(\bar{x})=s/\sqrt n.
\end{equation}
\textbf{人话：}平均值概括中心，标准差概括单次测量波动，标准误概括这个平均值有多不稳。样本越多，独立随机波动越容易被平均掉。
\textbf{符号：}$x_i$ 为第 $i$ 个观测；$n>1$ 为样本数；$\bar{x}$ 为样本均值；$s^2,s$ 为样本方差与标准差；$\operatorname{SE}$ 为标准误。除方差单位平方外，其他统计量与数据同单位。
\textbf{条件与误用：}上述标准误要求独立同分布观测及有限方差；$n-1$ 是估计均值后损失一个自由度产生的修正。每秒记录的相邻温度有相关性，不能把一万个点当一万次独立实验。基础阅读：\url{https://www.itl.nist.gov/div898/handbook/eda/section3/eda356.htm}。

\subsection{S12：大数定律——多测能改善什么}
\begin{equation}
 \bar X_n=\frac1n\sum_{i=1}^{n}X_i\xrightarrow{P}\mu.
\end{equation}
\textbf{人话：}独立重复次数越来越多，样本平均偏离真实平均很多的概率会越来越小。
\textbf{符号：}$X_i$ 是第 $i$ 次随机观测，$n$ 是次数，$\bar X_n$ 为随机的样本均值，$\mu=\mathrm E[X_i]$ 为总体均值；$\xrightarrow{P}$ 表示依概率收敛。均值与观测同单位。
\textbf{条件与误用：}这里采用独立同分布且绝对一阶矩有限的版本。它不保证每增加一个样本就更接近真值，也消不掉传感器整体偏高的系统误差。基础阅读：MIT 概率课程，\url{https://ocw.mit.edu/courses/6-041sc-probabilistic-systems-analysis-and-applied-probability-fall-2013/}。

\subsection{S13：中心极限定理——为什么均值常近似正态}
\begin{equation}
 \frac{\sqrt n(\bar X_n-\mu)}{\sigma}\xrightarrow{d}\mathcal N(0,1).
\end{equation}
\textbf{人话：}很多独立小波动的平均，在合适缩放后常接近钟形分布。注意变得近似正态的是样本平均的抽样分布。
\textbf{符号：}$n,\bar X_n,\mu$ 同上一节；$\sigma>0$ 为总体标准差；$\xrightarrow d$ 为分布收敛；$\mathcal N(0,1)$ 为均值零、方差一的正态分布。分式无单位。
\textbf{条件与误用：}此经典版本要求独立同分布、有限且非零方差。不存在适合所有数据的“超过 30 个就一定正态”规则；极偏、重尾或强相关数据可能需要很多样本或不同理论。基础阅读：\url{https://ocw.mit.edu/courses/6-041sc-probabilistic-systems-analysis-and-applied-probability-fall-2013/}。

\subsection{S14：总体均值的 $t$ 置信区间}
\begin{equation}
 \bar{x}\pm t_{1-\alpha/2,\,n-1}\frac{s}{\sqrt n}.
\end{equation}
\textbf{人话：}用“估计中心加减一段误差范围”报告平均值。假如反复采样并按此规则做区间，大约 $100(1-\alpha)\%$ 的区间会覆盖真实均值。
\textbf{符号：}$\bar x,s,n$ 是样本均值、样本标准差、样本数；$\alpha$ 是不覆盖概率的名义水平，常取 $0.05$；$t_{q,\nu}$ 是自由度 $\nu$ 的 $t$ 分布的 $q$ 分位数，即左侧累计概率为 $q$ 的位置；此处 $q=1-\alpha/2,\nu=n-1$。区间单位与数据一致。
\textbf{条件与误用：}独立正态样本时是精确区间，其他分布大样本时可能近似有效。它不是“95\% 的单个观测位于此区间”，也不是频率学框架下已固定真值在这个已算区间里的概率为 95\%。阅读：\url{https://www.itl.nist.gov/div898/handbook/prc/section1/prc14.htm}。

\subsection{S15：单样本 $t$ 检验}
\begin{equation}
 t_{\rm obs}=\frac{\bar{x}-\mu_0}{s/\sqrt n},\qquad
 p=2\bigl[1-F_{t,n-1}(|t_{\rm obs}|)\bigr].
\end{equation}
\textbf{人话：}检验观测平均与宣称值之间的差异，是否大到难以仅由随机波动解释。
\textbf{符号：}$\mu_0$ 是原假设声称的总体均值，与 $\bar x,s$ 同单位；$t_{\rm obs}$ 是算出的无量纲统计量；$F_{t,n-1}$ 是自由度 $n-1$ 的 $t$ 累积分布函数；竖线为绝对值；$p$ 是双侧 $p$ 值，无单位。
\textbf{条件与误用：}独立正态样本给出精确检验，备择假设为均值不等于 $\mu_0$。$p<\alpha$ 时拒绝原假设；$p\ge\alpha$ 不等于证明原假设正确。$p$ 不是“原假设为真的概率”。应同时报告差值大小、区间和实际意义。阅读：\url{https://docs.scipy.org/doc/scipy/reference/generated/scipy.stats.ttest_1samp.html}。

\subsection{S16：两独立样本的 Welch 检验}
\begin{equation}
 t=\frac{\bar{x}_1-\bar{x}_2-\Delta_0}{\sqrt{s_1^2/n_1+s_2^2/n_2}},\qquad
 \nu\approx\frac{(s_1^2/n_1+s_2^2/n_2)^2}{(s_1^2/n_1)^2/(n_1-1)+(s_2^2/n_2)^2/(n_2-1)}.
\end{equation}
\textbf{人话：}比较两种方案的平均表现，允许两组波动程度不同。
\textbf{符号：}下标 $1,2$ 代表两组；$\bar{x}_j,s_j^2,n_j$ 是第 $j$ 组均值、方差、数量；$\Delta_0$ 是原假设中的均值差，通常为零；$t$ 是统计量，$\nu$ 是近似自由度，二者无单位。两组数据必须表达同一量且统一单位。
\textbf{条件与误用：}两组相互独立、组内观测独立；正态或适当大样本近似。若同一台设备在改造前后各测一次，应优先对“成对差值”作单样本检验，不能抹掉配对结构。阅读：\url{https://docs.scipy.org/doc/scipy/reference/generated/scipy.stats.ttest_ind.html}。

\section{回归：从实验反推参数}
\subsection{S17：线性最小二乘}
\begin{equation}
 \mathbf y=X\boldsymbol\beta+\boldsymbol\varepsilon,\qquad
 \hat{\boldsymbol\beta}=\underset{\boldsymbol\beta}{\operatorname{argmin}}\,
 \|X\boldsymbol\beta-\mathbf y\|_2^2
 =(X^{\mathsf T}X)^{-1}X^{\mathsf T}\mathbf y.
\end{equation}
\textbf{人话：}找一组参数，让所有预测与观测的差的平方加起来最小。“线性”是对未知参数线性，可以包含已知的 $x^2$ 或 $\sin x$ 特征。
\textbf{符号：}$\mathbf y$ 是 $n$ 个响应组成的列向量；$X$ 是 $n\times p$ 设计矩阵，每行一个样本、每列一个已知特征；$\boldsymbol\beta$ 为 $p$ 个待估参数；$\boldsymbol\varepsilon$ 为观测误差；$\mathsf T$ 表示转置；$-1$ 表示矩阵逆；$\|\cdot\|_2^2$ 表示元素平方和；$\operatorname{argmin}$ 返回使目标最小的参数。第 $j$ 个参数单位为响应单位除以第 $j$ 列特征单位。
\textbf{条件与误用：}闭式逆矩阵表达要求满列秩，通常 $n\ge p$；误差条件均值为零才有无偏解释。实际程序用 QR、SVD 或最小二乘求解器，不显式求逆。高次多项式外推很容易失真。阅读：\url{https://www.itl.nist.gov/div898/handbook/pmd/section1/pmd141.htm}。

\subsection{S18：线性回归参数不确定性与新观测区间}
\begin{equation}
 \hat\sigma^2=\frac{\|\mathbf y-X\hat{\boldsymbol\beta}\|_2^2}{n-p},\qquad
 \widehat{\operatorname{Cov}}(\hat{\boldsymbol\beta})=\hat\sigma^2(X^{\mathsf T}X)^{-1},
\end{equation}
\begin{equation}
 \hat y_0\pm t_{1-\alpha/2,n-p}\hat\sigma\sqrt{1+\mathbf x_0^{\mathsf T}(X^{\mathsf T}X)^{-1}\mathbf x_0},\qquad
 \hat y_0=\mathbf x_0^{\mathsf T}\hat{\boldsymbol\beta}.
\end{equation}
\textbf{人话：}第一式估计参数波动；第二式预测未来一次实际观测的范围。根号内的 $1$ 表示新观测本身的随机噪声，去掉它才是均值响应的置信区间。
\textbf{符号：}$X,\mathbf y,\hat{\boldsymbol\beta},n,p$ 同上一节，$p$ 含截距参数；$\hat\sigma^2$ 为噪声方差估计；$\mathbf x_0$ 是新位置的 $p$ 维特征列；$\hat y_0$ 为预测中心；$\alpha$ 为显著性水平；$t$ 为对应分位数。协方差矩阵第 $i,j$ 项单位为参数 $i,j$ 单位的乘积；区间与响应同单位。
\textbf{条件与误用：}要求 $n>p$、满列秩、独立同方差正态误差，以及新观测独立且同噪声方差。不同方差或时间相关误差需修正协方差，不能照搬。阅读：\url{https://www.itl.nist.gov/div898/handbook/pmd/section1/pmd131.htm}。

\subsection{S19：非线性与加权最小二乘}
\begin{equation}
 \hat{\boldsymbol\theta}=\underset{\boldsymbol\theta\in\Omega}{\operatorname{argmin}}
 \sum_{i=1}^{n}\left[\frac{f(\mathbf x_i;\boldsymbol\theta)-y_i}{\sigma_i}\right]^2.
\end{equation}
\textbf{人话：}用物理模型 $f$ 拟合数据；测量越不准，给它的权重越小。例如拟合冷却曲线中的换热系数。
\textbf{符号：}$\boldsymbol\theta$ 为待估参数向量；$\Omega$ 是参数允许范围；$\mathbf x_i$ 为第 $i$ 组输入，可包含时间和环境温度；$y_i$ 为对应观测；$\sigma_i>0$ 为该观测误差的标准差；$n$ 为数据量；$i$ 为样本下标。$f,y_i,\sigma_i$ 单位相同，每项残差标准化后无单位。参数单位由模型决定。
\textbf{条件与误用：}独立高斯误差且标准差已知时，此目标对应最大似然；有相关误差应使用完整协方差矩阵。非线性问题可能有局部极小，需要合理边界、多初值和残差检查。参数强相关时，即使曲线拟合漂亮，也可能识别不出各个参数。阅读：\url{https://docs.scipy.org/doc/scipy/reference/generated/scipy.optimize.curve_fit.html}。

\subsection{S20：残差、RMSE 与决定系数}
\begin{equation}
 e_i=y_i-\hat y_i,\qquad
 \operatorname{RMSE}=\sqrt{\frac1n\sum_{i=1}^{n}e_i^2},\qquad
 R^2=1-\frac{\sum_{i=1}^{n}e_i^2}{\sum_{i=1}^{n}(y_i-\bar y)^2}.
\end{equation}
\textbf{人话：}残差是每一点算错多少；RMSE 汇总典型误差；$R^2$ 比较模型与直接用均值预测谁更好。
\textbf{符号：}$y_i$ 为观测，$\hat y_i$ 为预测，$e_i$ 为残差，$\bar y$ 为当前评价样本均值，$n$ 为评价样本量。残差和 RMSE 与响应同单位，$R^2$ 无单位。
\textbf{条件与误用：}$R^2$ 分母非零才有定义；测试集上的 $R^2$ 可以为负。必须看残差对时间、输入、预测值的散点图：趋势提示漏掉结构，喇叭口提示异方差，连续同号提示相关。训练集分数不能替代独立测试，时间数据按时间划分，避免未来信息泄漏。阅读：\url{https://scikit-learn.org/stable/modules/model_evaluation.html#regression-metrics}。

\section{让不确定性穿过物理模型}
\subsection{S21：一阶误差传播}
\begin{equation}
 u_y^2\approx\mathbf g^{\mathsf T}\Sigma\mathbf g
 =\sum_{i=1}^{d}\sum_{j=1}^{d}
 \frac{\partial f}{\partial x_i}\frac{\partial f}{\partial x_j}\Sigma_{ij},\qquad
 \mathbf g=\nabla f(\boldsymbol\mu).
\end{equation}
\textbf{人话：}输入有小误差时，先看输出对每个输入有多敏感，再把输入的不确定性传到输出。
\textbf{符号：}$y=f(\mathbf x)$ 为输出；$\mathbf x$ 有 $d$ 个输入分量；$\boldsymbol\mu$ 为输入均值向量；$\mathbf g$ 为在均值处的梯度，第 $i$ 项为偏导数；$\Sigma_{ij}=\operatorname{Cov}(X_i,X_j)$ 为输入协方差；$u_y$ 为输出标准不确定度；$i,j$ 枚举输入。偏导单位为输出单位除以对应输入单位，结果 $u_y^2$ 为输出单位平方。
\textbf{条件与误用：}需要均值附近可微、波动足够小使线性化合理。输入相关时不能删掉交叉项；强非线性、大范围波动或阈值问题宜做蒙特卡洛。这里的标准不确定度不是已知的真实误差。阅读：\url{https://doi.org/10.59161/JCGM100-2008E}。

\subsection{S22：乘幂模型的相对不确定性}
\begin{equation}
 y=C\prod_{i=1}^{d}x_i^{a_i},\qquad
 \left(\frac{u_y}{|y|}\right)^2\approx
 \sum_{i=1}^{d}a_i^2\left(\frac{u_i}{|x_i|}\right)^2.
\end{equation}
\textbf{人话：}乘除和乘幂模型中，独立输入的相对不确定性可以按指数加权后平方求和。例如面积随长度平方变化，同样的相对长度误差会被放大。
\textbf{符号：}$C$ 为当作已知的常数；$\prod$ 为连乘；$x_i\ne0$ 为输入估计值；$a_i$ 为已知指数；$u_i$ 为输入标准不确定度；$u_y$ 为输出标准不确定度；$d$ 为输入数。指数无单位，相对不确定度无单位，$C$ 的单位需使等式量纲一致。
\textbf{条件与误用：}这是上一条在输入互不相关且小误差下的特例，非整数幂通常要求正输入。$C$ 若也不确定，应加入传播；输入接近零时相对误差会失去实用意义。阅读：\url{https://doi.org/10.59161/JCGM100-2008E}。

\subsection{S23：蒙特卡洛估计与抽样误差}
\begin{equation}
 y^{(b)}=f(\mathbf x^{(b)}),\quad
 \hat\mu_y=\frac1B\sum_{b=1}^{B}y^{(b)},\quad
 s_y^2=\frac1{B-1}\sum_{b=1}^{B}(y^{(b)}-\hat\mu_y)^2,\quad
 \widehat{\operatorname{SE}}(\hat\mu_y)=s_y/\sqrt B.
\end{equation}
\textbf{人话：}按输入的联合分布抽很多组参数，每组都运行模型，收集输出分布。最后一个数只衡量模拟平均值因抽样次数有限而产生的误差。
\textbf{符号：}$b$ 为第几次模拟，括号上标不是乘方；$B>1$ 为模拟次数；$\mathbf x^{(b)}$ 为输入向量样本；$y^{(b)}$ 为对应输出；$\hat\mu_y$ 为模拟均值；$s_y$ 为输出标准差。所有输出统计量与 $y$ 同单位，$s_y^2$ 单位平方。
\textbf{条件与误用：}标准误表达适用于独立抽样和有限方差。必须保留输入相关性，并固定随机种子便于复现。$s_y$ 是输入不确定性导致的输出散布，不会因 $B$ 增大而消失；$s_y/\sqrt B$ 会减小。阅读：\url{https://doi.org/10.59161/JCGM101-2008}。

\subsection{S24：Bootstrap——让数据自己模拟重复实验}
\begin{equation}
 \hat\theta^{*(b)}=T(x_1^{*(b)},\ldots,x_n^{*(b)}),\qquad
 I_{\rm perc}=\left[Q^*_{\alpha/2},Q^*_{1-\alpha/2}\right].
\end{equation}
\textbf{人话：}从现有 $n$ 个数据中有放回地再抽 $n$ 个，计算一次想研究的统计量，重复很多次。星号统计量的分布近似估计量在重复采样中的波动。
\textbf{符号：}$T$ 是统计量计算规则，如均值或拟合参数；$x_i^{*(b)}$ 是第 $b$ 次重抽样的第 $i$ 个数据；$\hat\theta^{*(b)}$ 是对应估计；$Q_q^*$ 为所有重抽样估计的 $q$ 分位数；$\alpha$ 为名义不覆盖水平；$I_{\rm perc}$ 为百分位区间。区间单位与统计量一致。
\textbf{条件与误用：}普通独立重抽样要求原数据可视为独立同分布样本。配对数据必须成对抽，时间序列需考虑块 Bootstrap。百分位法便于理解但可能覆盖不足，偏差明显时考虑 BCa；极少数据或未观察过的尾部风险不能凭空补出来。阅读：\url{https://docs.scipy.org/doc/scipy/reference/generated/scipy.stats.bootstrap.html}。

\section{敏感性：哪个参数最值得认真测量}
\subsection{S25：局部敏感度与无量纲弹性}
\begin{equation}
 G_i=\frac{\partial f}{\partial x_i},\qquad
 E_i=\frac{x_i}{y}\frac{\partial f}{\partial x_i},\qquad
 G_i\approx\frac{f(\mathbf x+h\mathbf e_i)-f(\mathbf x-h\mathbf e_i)}{2h}.
\end{equation}
\textbf{人话：}只轻轻改变一个参数，观察结果如何变化。$E_i=2$ 表示在当前点附近，输入增大约 1\% 时输出约增大 2\%。
\textbf{符号：}$y=f(\mathbf x)$ 为输出；$i$ 为待研究参数编号；$G_i$ 为有量纲敏感度，单位是输出除以输入 $i$ 的单位；$E_i$ 为无量纲弹性；$\mathbf e_i$ 为第 $i$ 个坐标为一、其余为零的方向向量；$h$ 为第 $i$ 个输入的微小步长，与其同单位。
\textbf{条件与误用：}要求局部可微；弹性还要求 $y\ne0$。有限差分的 $h$ 既不能大到跨越非线性区域，也不能小到被数值噪声吞掉；应比较不同步长。这是某个基准点的结论，不是整个可行范围的全局重要性。阅读：\url{https://docs.scipy.org/doc/scipy/reference/differentiate.html}。

\subsection{S26：Sobol 一阶与总效应指标}
\begin{equation}
 S_i=\frac{\operatorname{Var}_{X_i}(\mathrm E[Y\mid X_i])}{\operatorname{Var}(Y)},\qquad
 S_{T_i}=\frac{\mathrm E_{X_{-i}}[\operatorname{Var}(Y\mid X_{-i})]}{\operatorname{Var}(Y)}.
\end{equation}
\textbf{人话：}一阶指标衡量参数自己造成多少输出波动；总效应还包含它与其他参数配合产生的波动。总效应大、一阶小，说明它主要通过交互作用影响结果。
\textbf{符号：}$Y=f(X_1,\ldots,X_d)$ 是随机输出；$X_i$ 是第 $i$ 个随机输入；$X_{-i}$ 指除第 $i$ 个以外的所有输入；$\mathrm E[\cdot\mid\cdot]$ 为条件期望；$\operatorname{Var}(\cdot\mid\cdot)$ 为条件方差；外侧下标说明对哪个输入的分布再平均或计算方差；$S_i,S_{T_i}$ 均无单位。
\textbf{条件与误用：}经典方差分解解释要求输入相互独立、输出有限且非零方差；分布与范围也是问题设定的一部分。改变范围可能改变排名。有限样本估计可能略为负或超过一，应看置信区间和收敛性，不能解释成负的真实贡献。Sobol 采样和分析必须配套，不能把随意乱序的实验表直接塞给估计器。阅读：\url{https://salib.readthedocs.io/en/stable/api/SALib.analyze.html#module-SALib.analyze.sobol}。

\section{A 题中一套可直接照做的统计检查顺序}
先把每一列的物理含义、单位、采样时刻、设备编号写清楚，保留原始数据。再画原始曲线，确认异常点是传感器错误还是模型必须解释的真实现象。选物理模型后，分出独立验证工况；只在训练数据上估参数，再对验证数据算误差。最后改变网格、求解容差、关键输入与分布假设，分别检查数值误差和参数不确定性。
报告结果时，把点估计、单位、区间、区间的含义、数据量和关键假设一起写。例如“某系数估计为……，95\% 置信区间为……，由独立重复实验并在正态误差假设下计算”。不要用一串小数位代替精度证据，也不要让显著性检验掩盖效应大小。


\chapter{基础定理与代表性前沿方法：读论文时抓住什么}
本章选取连接传统物理建模与科学机器学习的代表性方法。这里的“前沿”表示研究方向，不表示它们都是最近发表，更不是截至 2026 年的完整最新清单。竞赛中是否使用一种方法，取决于数据量、算力、可解释性和验证结果。先建立可复现的守恒方程与数值基线，再决定新方法能否降低误差或节省总计算成本。

\section{F01：弱形式——把难求的高阶导数转移出去}
对一维稳态导热，考虑
\begin{equation}
 \begin{aligned}
 -\frac{\mathrm d}{\mathrm dx}\!\left(k\frac{\mathrm dT}{\mathrm dx}\right)&=q,\\
 \int_0^L kT'v'\,\mathrm dx&=\int_0^L qv\,\mathrm dx.
 \end{aligned}
\end{equation}
弱式对所有满足 $v(0)=v(L)=0$ 的适当测试函数 $v$ 成立。
\textbf{人话：}原式要求每一点满足热量平衡；弱式先乘一个用于“检查误差”的函数，再对整段杆积分。分部积分把温度的二阶导数变成一阶导数，因此可以使用分段直线来逼近温度。
\textbf{符号与单位：}$x$ 是位置，$L$ 是杆长，单位 m；$T$ 是温度，单位 K；$k>0$ 是导热系数，单位 $\mathrm{W/(m\,K)}$；$q$ 是体积热源，单位 $\mathrm{W/m^3}$；$T'=\mathrm dT/\mathrm dx$；$v$ 是人为选择的无量纲测试函数，$v'=\mathrm dv/\mathrm dx$，单位 $\mathrm m^{-1}$；$\mathrm dx$ 是积分长度元。两边弱式单位均为 $\mathrm{W/m^2}$，这里采用恒截面积、单位截面积的写法。
\textbf{条件与误用：}本式取两端温度已知的 Dirichlet 边界，故测试函数端点为零，边界项才消失。绝热、给定热流或对流边界需要重新保留并处理边界项。弱解并不要求处处存在经典二阶导数，但仍要求适当可积性。它是有限元和弱形式辨识的基础，不能理解为“物理约束变松了可以任意不满足”。
\textbf{阅读定位：}Dokken 的 FEniCSx 教程给出 Poisson 方程的强式到弱式推导；上面是面向导热问题的原创改写。\url{https://jsdokken.com/dolfinx-tutorial/chapter1/fundamentals.html}。

\section{F02：SVD 与 Eckart--Young 最佳低秩逼近定理}
\begin{equation}
 A=U\Sigma V^{\mathsf T},\qquad
 A_r=\sum_{j=1}^{r}\sigma_j\mathbf u_j\mathbf v_j^{\mathsf T},\qquad
 \min_{\operatorname{rank}(B)\le r}\|A-B\|_F^2
 =\|A-A_r\|_F^2=\sum_{j=r+1}^{q}\sigma_j^2.
\end{equation}
\textbf{人话：}一大张时空温度表可能主要由少数几种变化模式组成。SVD 把它拆成按重要程度排列的模式，只留前几个时，在平方误差意义下已经是同等秩预算里最好的近似。
\textbf{符号与单位：}$A$ 是 $m\times n$ 数据矩阵，$m$ 为位置数、$n$ 为快照数；$q=\min(m,n)$；$U,V$ 的列 $\mathbf u_j,\mathbf v_j$ 是单位长度且相互正交的方向向量，无量纲；$\Sigma$ 为含奇异值的对角矩阵；$\sigma_1\ge\cdots\ge\sigma_q\ge0$ 是奇异值，单位与 $A$ 元素相同；$r$ 是保留模式数；$B$ 为任意竞争近似矩阵；$\operatorname{rank}$ 是独立方向数量；$\|\cdot\|_F^2$ 是所有元素平方和；$\mathsf T$ 是转置。
\textbf{条件与误用：}这是任意实矩阵在 Frobenius 范数下的有限维定理。若各列混合温度、速度、压力，必须先有合理尺度或加权内积，否则“最重要”主要由数值单位决定。低训练重构误差不自动保证新工况预测准确；奇异值相同时最佳解可能不唯一。
\textbf{阅读定位：}MIT 的矩阵方法课程第 7 讲专门讲解这一基础定理。\url{https://ocw.mit.edu/courses/18-065-matrix-methods-in-data-analysis-signal-processing-and-machine-learning-spring-2018/resources/lecture-7-eckart-young-the-closest-rank-k-matrix-to-a/}。

\section{F03：POD--Galerkin——把高维物理模型压缩成小模型}
\begin{equation}
 \mathbf u(t)\approx\bar{\mathbf u}+\Phi_r\mathbf a(t),\qquad
 \dot{\mathbf a}(t)=\Phi_r^{\mathsf T}F(\bar{\mathbf u}+\Phi_r\mathbf a(t),t),\qquad
 \eta_r=\frac{\sum_{j=1}^{r}\sigma_j^2}{\sum_{j=1}^{q}\sigma_j^2}.
\end{equation}
\textbf{人话：}把整条温度曲线表示成“平均曲线加上少数形状的加权组合”，只追踪这些权重如何随时间变化。
\textbf{符号与单位：}$\mathbf u(t)$ 是 $m$ 维离散场，$t$ 为时间 s；$\bar{\mathbf u}$ 是固定平均场；$\Phi_r$ 是由去均值快照 SVD 得到的前 $r$ 个左奇异向量组成的 $m\times r$ 矩阵，无量纲；$\mathbf a$ 是 $r$ 个模式幅值，与场同单位；$F$ 是原系统 $\dot{\mathbf u}=F(\mathbf u,t)$ 的右端，单位为场单位每秒；点号是时间导数；$\sigma_j,q$ 为去均值快照的奇异值及总数；$\eta_r$ 为无量纲快照平方幅值保留比例，分母须非零。
\textbf{条件与误用：}此简式假定欧氏正交基、原系统质量矩阵为单位矩阵、平均场不随时间变化。有限元带质量矩阵时需保留投影质量矩阵。保留 99\% 快照“能量”不等于任意物理量都误差小于 1\%；这里的能量通常仅指平方范数。对流、突变、新参数区间可能需要更多模式；非线性右端仍可能昂贵，不能只凭状态维数小就宣布计算提速。
\textbf{阅读定位：}Gubisch 与 Volkwein，2017，POD 基础与降阶。\url{https://doi.org/10.1137/1.9781611974829.ch1}。

\section{F04：PINN——把方程残差加入训练目标}
\begin{equation}
 r_{\boldsymbol\theta}=\frac{\partial u_{\boldsymbol\theta}}{\partial t}
 -a\frac{\partial^2u_{\boldsymbol\theta}}{\partial x^2}-s,\qquad
 \mathcal L=\frac{w_r}{N_r}\sum_{i=1}^{N_r}r_{\boldsymbol\theta}(x_i,t_i)^2
 +\frac{w_d}{N_d}\sum_{j=1}^{N_d}\left[u_{\boldsymbol\theta}(\tilde x_j,\tilde t_j)-\tilde u_j\right]^2
 +w_b\mathcal L_b+w_0\mathcal L_0.
\end{equation}
\textbf{人话：}神经网络猜一个完整温度场；训练时既惩罚它与观测不符，也惩罚它不满足热方程、边界和初值。网络对坐标的导数通常用自动微分计算。
\textbf{符号与单位：}本节先将全部坐标、场和方程无量纲化。$u_{\boldsymbol\theta}$ 为网络输出；$\boldsymbol\theta$ 为网络权重；$x,t$ 是无量纲位置和时间；$a$ 为缩放后的扩散系数；$s$ 为无量纲源项；$r$ 为方程残差；$(x_i,t_i)$ 是 $N_r$ 个物理检查点；$(\tilde x_j,\tilde t_j,\tilde u_j)$ 为 $N_d$ 个数据点，波浪号表示观测；$\mathcal L_b,\mathcal L_0$ 为边界、初值的平均平方残差；$w_r,w_d,w_b,w_0>0$ 是权重；所有损失项无量纲。
\textbf{条件与误用：}网络需有足够光滑的激活函数，以支持所需阶数的导数。残差小不自动等于解误差小；采样遗漏、权重失衡、优化失败都会影响结果。反问题中可把 $a$ 也设为待估参数，但仍需检验可辨识性。对一维标准导热，有限差分通常是值得先做的基线。
\textbf{原始文献：}Raissi、Perdikaris 与 Karniadakis，2017 预印本。\url{https://arxiv.org/abs/1711.10561}。

\section{F05：SINDy——从候选物理项中挑出简洁方程}
\begin{equation}
 \dot X\approx\Theta(X)\Xi,\qquad
 \hat\Xi=\underset{\Xi}{\operatorname{argmin}}
 \left\{\frac12\|\dot X-\Theta(X)\Xi\|_F^2+\lambda\|\Xi\|_1\right\}.
\end{equation}
\textbf{人话：}先列出可能出现的常数项、一次项、平方项、交叉项等，再寻找只用少数项就能解释变化率的方程。这里用 $\ell_1$ 惩罚展示稀疏回归思想，原始 SINDy 也采用序贯阈值最小二乘等做法。
\textbf{符号与单位：}先对状态和时间无量纲化；$X$ 是 $n\times d$ 状态数据矩阵，$n$ 为时刻数、$d$ 为状态数；$\dot X$ 为对应导数；$\Theta(X)$ 是 $n\times p$ 候选函数矩阵，$p$ 为候选项数；$\Xi$ 是 $p\times d$ 系数矩阵；$\lambda\ge0$ 控制稀疏程度；$\|\cdot\|_F^2$ 为全部元素平方和，$\|\cdot\|_1$ 此处明确指全部元素绝对值之和。全部量按无量纲表达，恢复物理单位时须撤销缩放。
\textbf{条件与误用：}真实动力学需要能被候选库合理表示，数据应激发足够多的动态。直接对噪声数据差分会放大噪声；候选项强相关可能导致错误选择。发现一个拟合好的方程不等于证明它是唯一真实机理，需要跨初值、跨工况验证。
\textbf{原始文献：}Brunton、Proctor 与 Kutz，2016。\url{https://doi.org/10.1073/pnas.1517384113}。

\section{F06：Weak SINDy——先积分，再辨识}
\begin{equation}
 \dot x(t)=\sum_{j=1}^{p}\xi_j\theta_j(x(t)),\qquad
 -\int_{t_0}^{t_1}x(t)\psi'(t)\,\mathrm dt
 =\sum_{j=1}^{p}\xi_j\int_{t_0}^{t_1}\theta_j(x(t))\psi(t)\,\mathrm dt.
\end{equation}
\textbf{人话：}为了避免直接给有噪声的观测求导，给方程乘一个自己选的光滑函数并积分，再通过分部积分把导数转移给这个已知函数。
\textbf{符号与单位：}此处使用无量纲标量状态与时间；$x(t)$ 为状态；$t_0,t_1$ 为积分窗口端点；$p$ 为候选项数；$\theta_j$ 为第 $j$ 个已知候选函数；$\xi_j$ 为它的未知系数；$\psi$ 为光滑测试函数且 $\psi(t_0)=\psi(t_1)=0$；$\psi'$ 为其时间导数；$\mathrm dt$ 为时间积分元。若不无量纲化，各系数单位必须使每个候选项与 $\dot x$ 同单位。
\textbf{条件与误用：}测试函数端点为零才可删掉分部积分边界项。窗口太宽可能抹掉局部动态，太窄可能不足以抑制噪声；仍需处理积分离散误差。这里展示易懂的一维 ODE 原理，PDE 的 WSINDy 会对时间和空间导数作相应弱形式处理，不是直接复制同一个积分就够了。
\textbf{原始文献：}Messenger 与 Bortz，2020 预印本，PDE 的 Weak SINDy。\url{https://arxiv.org/abs/2007.02848}。

\section{F07：高斯过程回归——同时预测数值和模型内不确定性}
\begin{equation}
 C=K+\sigma_n^2I,\qquad
 \mu_*=\mathbf k_*^{\mathsf T}C^{-1}\mathbf y,\qquad
 s_*^2=k(\mathbf x_*,\mathbf x_*)-\mathbf k_*^{\mathsf T}C^{-1}\mathbf k_*.
\end{equation}
\textbf{人话：}把未知函数看成一群可能曲线；观测数据排除不合适的曲线，剩下的平均给预测，分散程度给条件不确定性。
\textbf{符号与单位：}这里假定先验均值为零，非零均值可先扣除再加回；$\mathbf x_i$ 为第 $i$ 个输入，$\mathbf x_*$ 为查询输入；$n$ 为观测数；$\mathbf y$ 为 $n$ 个观测组成的向量；$K_{ij}=k(\mathbf x_i,\mathbf x_j)$ 为核协方差矩阵；$\mathbf k_*$ 的第 $i$ 项为 $k(\mathbf x_i,\mathbf x_*)$；$\sigma_n^2$ 为独立同方差观测噪声方差；$I$ 为 $n$ 阶单位矩阵；$\mu_*$ 为潜在函数预测均值，$s_*^2$ 为潜在函数预测方差。核和各方差单位为响应单位平方，均值与响应同单位。
\textbf{条件与误用：}公式基于高斯过程先验与高斯观测噪声，并把核超参数视为给定。若预测新的带噪观测，其方差再加 $\sigma_n^2$。这个不确定性依赖核、噪声和数据假设，不是现实误差保证；实现应用线性求解而非显式逆矩阵。普通稠密实现的计算量随样本数约呈立方增长。
\textbf{基础文献：}Rasmussen 与 Williams，2006，第 2 章。\url{https://gaussianprocess.org/gpml/chapters/RW2.pdf}。

\section{F08：贝叶斯优化的期望改进 EI}
\begin{equation}
 z(\mathbf x)=\frac{f_{\rm best}-\mu(\mathbf x)}{s(\mathbf x)},\qquad
 \operatorname{EI}(\mathbf x)=
 [f_{\rm best}-\mu(\mathbf x)]\Phi(z)+s(\mathbf x)\phi(z),\qquad
 \mathbf x_{\rm next}=\underset{\mathbf x\in\Omega}{\operatorname{argmax}}\operatorname{EI}(\mathbf x).
\end{equation}
\textbf{人话：}一个仿真很贵时，每次都要认真选下一个参数。EI 同时考虑“看起来能更好”和“这里还很不确定、值得探一探”。本式写的是最小化目标。
\textbf{符号与单位：}$\mathbf x$ 为设计参数向量；$\Omega$ 为可行参数域；$f_{\rm best}$ 为当前最小的无噪声目标观测；$\mu,s>0$ 为代理模型预测均值、标准差；$z$ 为无量纲标准化改进；$\Phi,\phi$ 为标准正态累积分布、密度；$\operatorname{EI}$ 为期望改进，与目标同单位；$\operatorname{argmax}$ 返回让 EI 最大的输入；$\mathbf x_{\rm next}$ 为下一次真正运行昂贵仿真的参数。
\textbf{条件与误用：}公式假定查询目标的后验为高斯分布；$s=0$ 时用极限 $\max(f_{\rm best}-\mu,0)$。强噪声下不能直接将偶然很低的观测当作确定最优，应采用考虑噪声的采集策略。它适合低到中等维度、单次评估昂贵的问题；便宜可微目标往往可直接优化。
\textbf{代表文献：}Snoek、Larochelle 与 Adams，2012。\url{https://arxiv.org/abs/1206.2944}。

\section{F09：DeepONet——学习“输入函数到输出函数”的映射}
\begin{equation}
 \mathcal G_{\boldsymbol\theta}(a)(x)
 =\sum_{j=1}^{p}b_j\bigl(a(\xi_1),\ldots,a(\xi_m)\bigr)\,t_j(x)+b_0.
\end{equation}
\textbf{人话：}普通回归常把几个数变成一个数；算子学习把一整条初温曲线变成另一条温度曲线，或者把材料系数场变成解场。分支网络理解输入函数，主干网络理解“想在什么位置看答案”。
\textbf{符号与单位：}本式采用无量纲化输入输出；$a$ 是输入函数；$\xi_1,\ldots,\xi_m$ 是 $m$ 个输入传感位置；$a(\xi_i)$ 为该处函数值；$x$ 是输出查询坐标，也可包含时间；$b_j$ 为分支网络第 $j$ 个输出；$t_j$ 为主干网络第 $j$ 个输出；$p$ 为共同特征数；$b_0$ 为偏置；$\boldsymbol\theta$ 为所有可训练参数；$\mathcal G_{\boldsymbol\theta}$ 为学到的算子。
\textbf{条件与误用：}需要多组“输入函数、对应解函数”的训练对。一条轨迹上的很多点不等于很多种工况。算子通用逼近结果在相应连续性、紧性和网络容量条件下说明存在足够好的表示，不保证梯度下降能找到它，也不保证训练范围之外准确。输入传感点分辨不了的细节可能无法恢复。
\textbf{代表文献：}Lu、Jin 与 Karniadakis，2019 预印本。\url{https://arxiv.org/abs/1910.03193}。

\section{F10：FNO——在频域学习非局部作用}
\begin{equation}
 v_{\ell+1}(x)=\sigma\!\left[
 W_\ell v_\ell(x)+\mathcal F^{-1}\!\left(R_\ell(k)\,\mathcal F[v_\ell](k)\right)(x)+b_\ell(x)
 \right].
\end{equation}
\textbf{人话：}先把场分解成不同空间频率，在频域让各通道相互混合，再变回空间。这使一个位置可以吸收远处的信息，适合学习某一类 PDE 的解算子。
\textbf{符号与单位：}本式采用无量纲坐标与隐特征；$\ell$ 为网络层编号；$v_\ell(x)$ 为该层在位置 $x$ 的多通道特征；$\sigma$ 为逐分量非线性激活函数，与奇异值符号无关；$W_\ell$ 为逐点通道混合矩阵；$\mathcal F$ 为傅里叶变换，$\mathcal F^{-1}$ 为逆变换；$k$ 是频率模式索引，不是导热系数；$R_\ell(k)$ 为保留频率上的可训练复矩阵；$b_\ell$ 为偏置。实际离散实现通常用 FFT，并只保留部分低频模式。
\textbf{条件与误用：}需准备覆盖目标问题族的训练数据；FFT 网格、周期性处理、填充和边界方式必须明确。更换分辨率可以复用一部分算子结构，但并不自动保证任意网格或任意高频细节正确。应与传统求解器比较总成本：数据生成加训练加预测，而非只比较训练完成后的单次推理。
\textbf{原始文献：}Li 等，2020 预印本、2021 ICLR 工作。\url{https://arxiv.org/abs/2010.08895}。

\section{怎样把文献中的公式写进自己的建模论文}
先说明你要解决的具体困难，例如未知参数、噪声导数、重复仿真昂贵或空间维数过大。随后写出采用公式，并重新定义本文符号、单位、边界条件和离散方式。最后写清它在自己的数据上通过了什么验证。引用作者和方法名称不能代替推导；公式外形相同也不代表物理假设相同。
学习顺序建议为：守恒方程与数值收敛，最小二乘与统计不确定性，弱形式与 SVD，最后再选与问题相符的科学机器学习方向。上述十节中，F01--F03 是重要基础，F04--F10 是连接这些基础的代表性研究方法。所有解释均为针对本手册的教学性重述，未逐字搬录论文，也未声称复现论文全部理论与实验。

\chapter{Python 入门与完整逐行注释模板}
\section{先看懂程序的几个基本规则}
\texttt{import} 是加载工具箱；\texttt{as np} 是为 NumPy 取简称；\texttt{def} 定义可反复使用的函数；冒号后的缩进表示函数或循环的内部。\texttt{return} 把计算结果交给调用者。Python 区分大小写。

单个等号是赋值；两个等号是判断相等。\texttt{**2} 是平方；\texttt{*} 是乘法；\texttt{/} 是除法。数学中的并列相乘在代码里必须写乘号。\texttt{\#} 后的文字是给人看的解释，不参与运算。\texttt{assert} 后面的条件不成立时就停止程序，避免把失败结果当答案。

数组下标从零开始，\texttt{y[0]} 表示第一个分量，\texttt{y[-1]} 表示最后一个。切片 \texttt{u[1:-1]} 表示去掉两端；\texttt{u[:-2]} 去掉最后两个；\texttt{u[2:]} 去掉最前两个，它们可分别与内部点的左邻居、右邻居对应。

\texttt{np.array} 创建数值数组；\texttt{np.linspace(a,b,n)} 在含端点的区间生成 $n$ 个等距点，其中 $a,b$ 为端点、$n$ 为点数；\texttt{np.exp} 与 \texttt{np.log} 是指数与自然对数；\texttt{np.abs} 是绝对值；\texttt{np.max} 是最大值；\texttt{np.all} 要求全部元素都满足条件。\texttt{lambda z: ...} 是短函数写法，冒号后是返回的表达式。

\section{运行步骤}
安装 Python 3.11 或以上版本后，在本手册目录打开终端。先运行安装命令，再运行相应文件。命令不是在 Python 的三个大于号交互提示符后输入，而是在系统终端输入。下面每行的井号后都解释该行目的。
\begin{Verbatim}[breaklines=true,breakanywhere=true,fontsize=\small]
python -m pip install numpy scipy  # 安装三个数值示例所需的数组与科学计算库
python code/ode_template.py  # 求解冷却ODE并核验解析解与事件时间
python code/heat_1d.py  # 比较三种热差分方法并检查网格加密误差
python code/optimize_template.py  # 运行带约束优化和冷却参数拟合
\end{Verbatim}
网络不可用时，需先在有网络的环境准备匹配 Python 版本的安装包。模板自身不访问网络，也不需要 API 密钥。输出中的误差允许末尾数字随平台略变；是否通过断言比完全逐字相同更重要。

\section{ODE：冷却到指定温度}
输入为环境温度、初始温度、冷却系数与目标温度；输出为误差和到达目标的时间。替换方程时同时修改状态初值、参数单位和事件函数。\texttt{rtol} 是相对误差控制，\texttt{atol} 是绝对误差控制；它们控制内部误差估计，不是最终误差的严格保证。
\begin{Verbatim}[breaklines=true,breakanywhere=true,fontsize=\small,numbers=left,numbersep=5pt]
import numpy as np  # 导入数组和指数函数工具，后文用 np 作为简称
from scipy.integrate import solve_ivp  # 导入自动选择时间步长的常微分方程求解器

ambient = 20.0  # 设置恒定环境温度，单位为摄氏度
initial = 90.0  # 设置物体初始温度，单位为摄氏度
rate = 0.08  # 设置冷却系数，单位为每秒
threshold = 40.0  # 设置停止计算的目标温度，单位为摄氏度

def rhs(t, y):  # 定义变化率函数，t 是秒，y 是保存各状态的数组
    return [-rate * (y[0] - ambient)]  # 返回温度变化率列表，单位为摄氏度每秒

def hit_threshold(t, y):  # 定义事件函数，求解器寻找这个函数等于零的时刻
    return y[0] - threshold  # 温度减去阈值，单位为摄氏度

hit_threshold.terminal = True  # 第一次满足事件条件时终止积分
hit_threshold.direction = -1  # 只检测温度从阈值上方向下穿过的事件
times = np.linspace(0.0, 60.0, 301)  # 生成希望输出的时间点，这不等于求解器内部步长
sol = solve_ivp(rhs, (0.0, 60.0), [initial], method="RK45", t_eval=times, events=hit_threshold, rtol=1e-8, atol=1e-10, max_step=0.5)  # 积分并限制内部最大步长为零点五秒
assert sol.success, sol.message  # 如果积分失败，立即显示求解器的错误说明
assert len(sol.t_events[0]) == 1, "未检测到目标温度事件"  # 验证恰好找到一个停止事件
exact = ambient + (initial - ambient) * np.exp(-rate * sol.t)  # 计算同一批输出时刻的解析温度
max_error = np.max(np.abs(sol.y[0] - exact))  # 求数值温度与解析温度的最大绝对误差，单位为摄氏度
event_exact = np.log((initial - ambient) / (threshold - ambient)) / rate  # 计算理论到达阈值时间，单位为秒
event_numeric = sol.t_events[0][0]  # 从事件结果读取数值到达时间，不能用最后一个输出采样点代替
assert max_error < 1e-5, "温度误差超出允许范围"  # 用独立解析解检查数值计算的正确性
assert abs(event_numeric - event_exact) < 1e-5, "事件时间误差超出允许范围"  # 检查事件定位精度
print("最大温度误差（摄氏度）:", max_error)  # 显示温度误差供论文验证使用
print("到达目标温度的时间（秒）:", event_numeric)  # 显示题目常要求的首次达到时刻
\end{Verbatim}

\section{一维热传导：显式、隐式与 Crank--Nicolson}
本模板只实现均匀材料、常热扩散率、无热源、两端恒为零的教学验证问题。温度也可理解为相对于固定端点温度的超温。它不是可直接承接任意对流边界、多层材料的通用成品；改变这些物理条件时须按正文重新装配系数与右端项。
\begin{Verbatim}[breaklines=true,breakanywhere=true,fontsize=\small,numbers=left,numbersep=5pt]
import numpy as np  # 导入数组、三角函数和指数函数
from scipy.linalg import solve_banded  # 导入带状矩阵求解器，避免对三对角矩阵求逆

length = 1.0  # 设置杆长，单位为米
alpha = 0.01  # 设置热扩散率，单位为平方米每秒
amplitude = 10.0  # 设置初始温度波形的最大幅值，单位为摄氏度
end_time = 0.5  # 设置最终计算时间，单位为秒

def run(intervals, method):  # 定义一次热传导计算，intervals 是空间区间数，method 是方法名称
    dx = length / intervals  # 用总长除以区间数得到空间步长，单位为米
    target_dt = 0.4 * dx**2 / alpha  # 选择满足显式稳定条件的目标时间步长，单位为秒
    steps = int(np.ceil(end_time / target_dt))  # 向上取整时间步数以保证实际步长不超过目标值
    dt = end_time / steps  # 调整实际步长使最后一步恰好到达终点
    ratio = alpha * dt / dx**2  # 计算无量纲扩散比，也叫网格傅里叶数
    assert ratio <= 0.5, "显式稳定性条件不满足"  # 检查本模板使用的统一时间网格满足显式条件
    x = np.linspace(0.0, length, intervals + 1)  # 生成含左右端点的空间坐标
    u = amplitude * np.sin(np.pi * x / length)  # 设置两端为零的正弦初温，单位为摄氏度
    u[[0, -1]] = 0.0  # 把浮点三角函数产生的端点舍入误差清零
    theta = {"explicit": 0.0, "implicit": 1.0, "cn": 0.5}[method]  # 统一表达显式、后向欧拉和克兰克尼科尔森方法
    count = intervals - 1  # 计算内部未知温度点的个数
    band = np.zeros((3, count))  # 创建三行带状矩阵存储格式，分别存上对角、主对角和下对角
    band[0, 1:] = -theta * ratio  # 填写上对角系数，第一个存储位置不使用
    band[1, :] = 1.0 + 2.0 * theta * ratio  # 填写主对角系数
    band[2, :-1] = -theta * ratio  # 填写下对角系数，最后一个存储位置不使用
    for step in range(steps):  # 按时间顺序重复更新温度场
        rhs = u[1:-1] + (1.0 - theta) * ratio * (u[:-2] - 2.0 * u[1:-1] + u[2:])  # 计算旧时刻温度形成的右端项
        u[1:-1] = solve_banded((1, 1), band, rhs)  # 解三对角方程并写入内部新温度，两端保持零
    exact = amplitude * np.sin(np.pi * x / length) * np.exp(-alpha * (np.pi / length)**2 * end_time)  # 计算精确的正弦衰减解
    error = np.max(np.abs(u - exact))  # 计算所有空间点的最大温度误差
    assert np.min(u) >= -1e-10, "出现不合理的负温度"  # 对本非负初边值例题检查非负性
    assert np.max(u) <= amplitude + 1e-10, "出现不合理的温度增长"  # 对本无热源例题检查最大值不会增加
    assert np.max(np.abs(u[[0, -1]])) < 1e-12, "边界条件未保持"  # 检查两端恒温条件
    return error  # 返回该网格的误差，单位为摄氏度

for method in ("explicit", "implicit", "cn"):  # 依次比较三个时间离散方法
    coarse_error = run(40, method)  # 用四十个空间区间计算较粗网格误差
    fine_error = run(80, method)  # 空间步长减半且时间步长约缩小到四分之一
    observed_order = np.log(coarse_error / fine_error) / np.log(2.0)  # 根据误差比例估计联合加密时的空间收敛阶
    assert fine_error < coarse_error / 3.0, "网格加密没有获得预期误差下降"  # 验证与二阶联合收敛相容的误差降低
    print(method, "粗网格误差=", coarse_error, "细网格误差=", fine_error, "联合收敛阶=", observed_order)  # 打印验证结果
\end{Verbatim}

\section{非线性优化与参数拟合}
第一个例子有手算最优解；第二个例子用已知参数生成无噪声数据，检查能否找回参数。实际有噪声数据不会精确恢复真值，还应增加置信区间、独立验证与可辨识性分析。
\begin{Verbatim}[breaklines=true,breakanywhere=true,fontsize=\small,numbers=left,numbersep=5pt]
import numpy as np  # 导入数组和范数计算工具
from scipy.optimize import minimize, least_squares  # 导入约束优化和残差向量拟合求解器

def objective(z):  # 定义无量纲目标函数，z 是两个决策变量组成的数组
    return (z[0] - 2.0)**2 + (z[1] - 1.0)**2  # 计算到理想点二一的平方距离，越小越好

def gradient(z):  # 定义目标函数的一阶偏导数组
    return np.array([2.0 * (z[0] - 2.0), 2.0 * (z[1] - 1.0)])  # 返回分别对两个变量求偏导的结果

constraint = {"type": "ineq", "fun": lambda z: 2.0 - z[0] - z[1], "jac": lambda z: np.array([-1.0, -1.0])}  # SciPy 的不等式函数必须大于等于零，表示两变量之和不超过二
starts = [np.array([0.2, 0.2]), np.array([1.8, 0.1]), np.array([0.1, 1.8])]  # 给出三个不同的可行初值供交叉验证
results = []  # 创建列表保存每次优化的完整结果
for start in starts:  # 从每个初值独立求解同一问题
    result = minimize(objective, start, jac=gradient, method="SLSQP", bounds=[(0.0, 2.0), (0.0, 2.0)], constraints=[constraint], options={"ftol": 1e-12, "maxiter": 200})  # 指定目标、梯度、边界、不等式和迭代精度
    assert result.success, result.message  # 如果优化器未正常终止就显示失败原因
    assert constraint["fun"](result.x) >= -1e-8, "不满足总量约束"  # 独立检查最终点是否可行
    assert np.all(result.x >= -1e-8), "不满足非负边界"  # 独立检查非负性
    results.append(result)  # 保存此次计算结果供比较
best = min(results, key=lambda result: result.fun)  # 从所有可行成功结果中选取目标值最小者
assert np.linalg.norm(best.x - np.array([1.5, 0.5])) < 1e-6, "未接近解析最优解"  # 用严格凸例题的解析解检验算法实现
print("约束最优变量:", best.x, "目标值:", best.fun)  # 输出最佳决策变量和目标值

times = np.linspace(0.0, 30.0, 31)  # 生成三十一个观测时间，单位为秒
data = 20.0 + 70.0 * np.exp(-0.08 * times)  # 生成无噪声教学数据，实际比赛应替换为真实温度测量值
sigma = 0.5  # 设置温度误差的参考标准差，单位为摄氏度

def residual(z):  # 定义残差向量函数，z 的两个分量都已经无量纲化
    cooling_rate = 0.1 * z[0]  # 把第一个无量纲参数换回每秒的冷却系数
    ambient_temperature = 20.0 * z[1]  # 把第二个无量纲参数换回摄氏度的环境温度
    prediction = ambient_temperature + (90.0 - ambient_temperature) * np.exp(-cooling_rate * times)  # 计算给定参数下的模型温度
    return (prediction - data) / sigma  # 返回按误差尺度归一化的无量纲残差数组

fit = least_squares(residual, [0.5, 0.5], bounds=([0.01, 0.0], [5.0, 3.0]), method="trf", xtol=1e-12, ftol=1e-12, gtol=1e-12)  # 在物理允许的参数范围内拟合残差向量
assert fit.success, fit.message  # 检查非线性最小二乘是否成功终止
estimated = np.array([0.1 * fit.x[0], 20.0 * fit.x[1]])  # 恢复有物理单位的冷却系数和环境温度
assert np.allclose(estimated, [0.08, 20.0], rtol=1e-6, atol=1e-6), "无噪声参数恢复失败"  # 用已知生成参数检查参数识别流程
print("拟合冷却系数（每秒）和环境温度（摄氏度）:", estimated)  # 输出拟合结果和对应单位
\end{Verbatim}

\chapter{论文落地与查阅路线}
\section{按问题找工具}
\begin{longtable}{p{0.24\textwidth}p{0.31\textwidth}p{0.35\textwidth}}
\toprule 题目现象 & 起步模型 & 必须核查 \\ \midrule
温度随时间变 & 集总热容 ODE & Biot 数、环境是否恒定 \\
物体内部温差明显 & 热传导 PDE 与差分 & 边界、热源、材料界面、收敛 \\
振动、轨迹、运动 & 牛顿方程与一阶系统 & 力的方向、坐标、接触事件 \\
管道输水、压差 & 连续方程与能量方程 & 雷诺数、损失、适用流态 \\
光强、成像、干涉 & 对应几何光学或波动光学 & 波长、角度、相干性与单位 \\
已知曲线想求参数 & 非线性最小二乘 & 误差结构、参数相关性 \\
选尺寸、功率、时间 & 带约束非线性优化 & 可行性、多初值、物理限制 \\
重复试验或测量噪声 & 概率模型与统计推断 & 独立性、分布假设、样本量 \\
数据多但机理不全 & SINDy、PINN 等选读 & 先有基线、数据覆盖与验证成本 \\
\bottomrule
\end{longtable}

\section{把公式写进论文的模板}
可以这样写：“将研究对象视为……，忽略……，依据……守恒建立……方程。式中各符号与单位列于符号表；初始条件由……获得，边界条件对应……。参数采用……数据估计。使用……方法求解，并通过解析特例、网格加密与独立观测检验。最终输出……，适用于……范围。”省略号必须由题目事实补足。

不要把代码运行成功写成“模型正确”；不要把一次优化终止写成“全局最优”；不要把文献中的某个材料参数直接搬到另一材料；不要把未经验证的前沿模型包装成必然优于经典方法。

\section{交付文件说明}
主文档是 PDF；\texttt{main.tex} 是可编辑 LaTeX 源文件；\texttt{chapters} 内保存章节；\texttt{code} 内是可独立运行的完整模板。公式均保留为 LaTeX 数学表达式，可以复制到论文。所有公式的符号含义以紧邻正文为准；无量纲比例和编号也已说明。

\section{基础文献的进一步阅读}
\begin{enumerate}
\item Boyd, S.; Vandenberghe, L. \emph{Convex Optimization}. Cambridge University Press, 2004. 凸性、对偶与 KKT。作者开放教材：\url{https://web.stanford.edu/~boyd/cvxbook/}。
\item NIST/SEMATECH. \emph{e-Handbook of Statistical Methods}. 统计分布、拟合、诊断与区间。\url{https://www.itl.nist.gov/div898/handbook/}。
\item SciPy 官方文档。实现时以实际安装版本为准，重点查 \texttt{solve\_ivp}、\texttt{solve\_banded}、\texttt{minimize}、\texttt{least\_squares}。\url{https://docs.scipy.org/doc/scipy/reference/}。
\item OpenStax. \emph{University Physics}, Volumes 1--3. 基础力学、流体、热学、光学。\url{https://openstax.org/subjects/science}。各物理章节另列具体入口。
\end{enumerate}
本文为教学性综合整理，文献观点用自己的语言概括；方法公式不是其实际性能的保证。网络来源核查日期为 2026 年 9 月 10 日；前沿章节是代表性路线，未声称穷尽或代表最新全部进展。

\end{document}
